\documentclass[letterpaper]{amsart}

%------ 
%package 
%------
\usepackage[utf8]{inputenc}
\usepackage[english]{babel}
\usepackage{amsmath,amssymb,amsthm,mathrsfs}
\usepackage{esint}
\usepackage[colorlinks,linktocpage]{hyperref}
\usepackage[all]{xy}
\usepackage{import}
\usepackage{pdfpages}
\usepackage{transparent}
\usepackage{xcolor}
\usepackage{comment}
\usepackage{upgreek}

%------
%theorem 
%------ 
\newtheorem{thm}{Theorem}[section]
\newtheorem{lem}[thm]{Lemma}
\newtheorem{cor}[thm]{Corollary}
\newtheorem{prop}[thm]{Proposition}
\newtheorem{conj}[thm]{Conjecture}

\theoremstyle{definition}
\newtheorem{defn}{Definition}[section]
\newtheorem{prob}{Problem}[section]
\newtheorem{exmp}{Example}[section]

\theoremstyle{remark}
\newtheorem{rmk}{Remark}[section]
\newtheorem*{rmk*}{Remark}
\newtheorem{fact}{Fact}[thm]
\newtheorem*{fact*}{Fact}
\newtheorem*{note}{Note}

%------
%rename
%------
\def\C{\mathbb{C}}
\def\R{\mathbb{R}}
\def\Q{\mathbb{Q}}
\def\Z{\mathbb{Z}}
\def\CP{\mathbb{CP}}
\def\supp{\mathrm{supp}}
\def\tr{\mathrm{tr}}
\def\diam{\mathrm{diam}}
\def\dim{\mathrm{dim}}
\def\Area{\mathrm{Area}}
\def\Ric{\mathrm{Ric}}
\def\Vol{\mathrm{Vol}}
\def\Rm{\mathrm{Rm}}
\def\Sph{\mathcal{S}}
\def\Hess{\mathrm{Hess}}
\def\Id{\mathrm{Id}}
\def\inj{\mathrm{inj}}
\def\Iso{\mathrm{Iso}}
\def\Diff{\mathrm{Diff}}
\def\Fix{\mathrm{Fix}}
\def\dV{\mathrm{dvol}}
\def\WL{\overset{\ast}{\rightharpoonup}}

\title{Einstein metrics with torus symmetry on $S^4$ and $\mathbb R^4$}
\author{Xiuxiong Chen}
\author{Conghan Dong}
%\address{Department of Mathematics, Duke University, 120 Science Dr, Durham, NC 27710, USA}
%\email{conghan.dong@duke.edu}
\begin{document}
\date{\today}

\begin{abstract}
This note gives direct proofs that, under an effective isometric $T^2$ action,
Einstein metrics on smooth homotopy four-spheres are round and complete
Ricci-flat metrics on smooth manifolds homeomorphic to $\mathbb R^4$ are
Euclidean or positive one-center Taub--NUT, up to isometry, scaling, and
orientation reversal. 
\end{abstract}

\maketitle

\section{Introduction}
A Riemannian metric $g$ is Einstein if $\mathrm{Ric}_g=\lambda g$ for
some constant $\lambda$. The classification of such metrics is a basic
problem in Riemannian geometry; see \cite{notes-Besse}. In dimension four,
even the simplest underlying manifolds lead to natural rigidity
questions. For $S^4$, the classical question is whether an Einstein
metric must be round, up to isometry and scaling. For $\mathbb R^4$,
the complete Ricci-flat examples include both the Euclidean metric and
the Taub--NUT metric \cite{notes-Hawking,notes-GibbonsHawking}.
In this note, we give proofs of the corresponding rigidity statements
under an effective isometric $T^2$ action.

\begin{thm}\label{thm-sphere}
Let $M$ be a smooth closed homotopy four-sphere, and let $g$ be an Einstein metric
on $M$. If $(M,g)$ admits an effective isometric $T^2$ action, then $(M,g)$
is isometric, after constant rescaling, to the unit round four-sphere.
\end{thm}

Liu proved the sphere rigidity statement under the additional assumption
of nonnegative sectional curvature, and established the polar reduction
for closed simply connected toric Einstein four-manifolds without that
curvature assumption \cite[Theorems~1.1--1.2]{notes-Liu}.
The sphere case is also contained in Li--Sun's uniqueness theorem
\cite[Theorem~1.5]{notes-LiSunCompact}; they record a proof by Broder and
Lafuente in \cite[Remark~D.19]{notes-LiSunCompact}.

For complete Ricci-flat metrics on manifolds homeomorphic to
$\mathbb R^4$, Chen--Li \cite[Theorem~1.6]{notes-ChenLi} obtained the
Euclidean or Taub--NUT classification under the curvature bound
\[
 |\mathrm{Rm}_g|(x)\leq\frac{k(r(x))}{r(x)^2},
 \qquad
 \int_1^\infty\frac{k(s)}{s}\,ds<\infty,
\]
where $r$ is distance from a fixed point and $k$ is positive and
nonincreasing, provided that the tangent cone at infinity is not
$\mathbb R\times[0,\infty)$. Their theorem does not assume symmetry.
Here we assume torus symmetry and impose no curvature-decay or
asymptotic hypothesis.

\begin{thm}\label{notes-R4-classification}
Let $(M^4,g)$ be a smooth complete Ricci-flat manifold homeomorphic to
$\mathbb R^4$. If $(M,g)$ admits an effective isometric $T^2$ action, then
$g$ is Euclidean or positive one-center Taub--NUT, up to isometry,
constant rescaling, and orientation reversal.
\end{thm}

The noncompact conclusion can also be recovered from Li--Sun's end
analysis and harmonic-map theory \cite{notes-LiSunCompact,notes-LiSunHarmonic},
together with the energy and curvature estimates of Weber
\cite{notes-Weber} and Cheeger--Tian \cite{notes-CheegerTian}.
Our approach gives an intrinsic proof through Green functions, sharp
harmonic-map estimates, and curvature transgression.

We briefly describe the two arguments. The torus-action classification
of Orlik--Raymond \cite{notes-OrlikRaymond} determines the orbit space
and the collapsing circle along each boundary arc. In the compact case,
the Einstein constant is positive, so we normalize
$\mathrm{Ric}_g=3g$. Write $h$ for the orbit metric and
$\rho=\sqrt{\det G}$ for the orbit-area function, where $G$ is the Gram
matrix of the two Killing fields. The reduced Einstein equation gives
$K_h\leq1$ and
\[
 \nabla_h^2(-\log\rho)\geq2h.
\]
Thus $-\log\rho$ is a strictly convex exhaustion of the interior of the
orbit space. We use this to show that its metric completion is
$\mathrm{CAT}(1)$. The two boundary arcs join the same fixed points and
are local geodesics. The $\mathrm{CAT}(1)$ property then forces the
diameter to be at least $\pi$. Bonnet--Myers and Cheng's maximal-diameter
theorem \cite{notes-Cheng} give roundness.

In the Ricci-flat case, $\rho$ is harmonic. Let $\tau$ be its harmonic
conjugate, normalized to vanish at the fixed point $p$. We first identify
the minimal positive Green function, normalized by
$-\Delta_g G_p=4\pi^2\delta_p$, as
\[
 G_p=\frac{1}{2\sqrt{\rho^2+\tau^2}}.
\]
Li--Yau's Green-function estimate \cite{notes-LiYau} then gives proper
global Weyl coordinates, and Colding's gradient estimate
\cite{notes-Colding} gives a sharp lower bound for the orbit metric.
The normalized matrix $H=\rho^{-1}G$ defines a harmonic map into
$\mathbb H^2$. A sharp asymptotic estimate for its $\tau$-derivative separates the
Euclidean case from a collapsing case. In the latter case, estimates for
suitable torus Killing fields control the boundary terms in the Euler
and Pontryagin transgression formulas. The Euler identity first gives
finite curvature energy. The Pontryagin identity then forces one Weyl
half to vanish, and simple connectivity gives a hyperk\"ahler structure.
A circle subgroup of the torus preserves the hyperk\"ahler triple and is
tri-Hamiltonian. Swann's complete classification
\cite[Theorem~3.1]{notes-Swann}, with the center description in
\cite[Section~2.2]{notes-DancerSwann}, then applies. The topology of
$\mathbb R^4$ forces a single center, leaving precisely the Euclidean
and Taub--NUT cases.

Section~\ref{sec-sphere} proves Theorem~\ref{thm-sphere}.
Section~\ref{sec-R4} develops the global coordinates, analytic estimates,
and curvature identities used to prove Theorem~\ref{notes-R4-classification}.

\subsection*{AI disclosure}
This note arose from an experiment in mathematical research through
interaction with ChatGPT. The author posed the problems and follow-up
questions, while ChatGPT developed the proposed mathematical arguments
and generated the initial manuscript. Its role was therefore substantial
and was not limited to language editing.
The author worked through the arguments and rewrote the main body in their own words, based on their understanding.

\section{\texorpdfstring{Einstein metrics on $S^4$}{Einstein metrics on S4}}\label{sec-sphere}
This section proves Theorem~\ref{thm-sphere}.

\subsection{The action and orthogonal transitivity}\label{subsec-orthogonal}
Write $\mathrm{Ric}_g=\lambda g$. The Einstein constant is positive.
Indeed, a nonzero Killing field $X$ generated by a circle subgroup satisfies $\frac12\Delta_g|X|_g^2=|\nabla^gX|_g^2-\mathrm{Ric}_g(X,X)$. Integrating gives
\begin{align*}
	\int_{M}|\nabla^g X|_g^2\,dV_g = \lambda \int_{M} |X|_g^2\,dV_g ,
\end{align*} 
which excludes $\lambda <0$, and the equality would make a nonzero parallel Killing field, contradicting $H^1(M; \mathbb{R}) =0$.

After rescaling, we can assume that $\mathrm{Ric}_g = 3 g$. Orlik--Raymond applies because $H_1(M;\mathbb Z)=0$: the quotient is a disk with no exceptional finite-isotropy orbits, and $\chi(M)=2$ gives two fixed points. Their smooth equivariant classification identifies the smoothing and action with the standard sphere \cite[\S\S4.2, 5.1--5.4]{notes-OrlikRaymond}. Thus, up to a diffeomorphism and a torus automorphism, we can assume that the effective $T^2$-action is given by the standard action:
\begin{align*}
	\left( e ^{i s_1}, e ^{i s_2} \right) \cdot  \left( z_1, z_2, t \right) = \left( e ^{i s_1} z_1, e ^{i s_2}z_2, t \right) ,
\end{align*} 
where $s_1,s_2\in\mathbb R/2\pi\mathbb Z$ are the torus action parameters and $M=S^4 = \{ (z_1, z_2, t) \in \mathbb{C}^2 \times \mathbb{R}: |z_1|^2+ |z_2|^2+ t ^2 =1\} $.

Let $B = M / T^2$ be the orbit space and $\pi : M \to B$ be the quotient map. Then $B$ is a closed $2$-disk, the action is free over the interior  $B ^{\circ}$ and the boundary $\partial B$ consists of two arcs $\Gamma _1, \Gamma _2$ joining two corners $P ^{\pm}$. More precisely, one can identify $\Gamma _1$ with the orbit of the invariant $2$-sphere $\{(z_1, 0, t): |z_1|^2 + t ^2= 1\}$, $\Gamma _2$ with the orbit of the invariant $2$-sphere $\{(0, z_2, t): |z_2|^2+ t ^2=1\}$, and $P ^{\pm}$ with the orbit of the fixed points $\{(0,0, \pm 1)\}$. 

Let $X, Y$ be the Killing fields of the two coordinate circles, and let $\alpha=g(X,\cdot)$ and $\beta=g(Y,\cdot)$ be their metric dual $1$-forms. In these coordinates, 
$$X(z_1, z_2, t) = (i z_1, 0, 0),\ \  Y(z_1, z_2, t) = (0, i z_2, 0).$$
The twist $1$-form associated with $X$ is 
$$\omega _{X} = * (\alpha \wedge d \alpha ),$$
where $*=*_g$ is the Hodge star operator.
Similarly, we have twist form $\omega _{Y}$.

For a Killing field, the twist identity says
\begin{align*}
	d \omega _{X} = 2 * \left( \alpha \wedge \mathrm{Ric}_g(X, \cdot ) \right) .
\end{align*} 
\begin{rmk}
	This twist identity also follows from the Hessian identity for a Killing field $X$:
\begin{align*}
	\nabla ^2_{V, W}X := \nabla _{V} \nabla _{W} X - \nabla _{\nabla _{V} W} X = R(V, X) W.
\end{align*} 
Taking the trace of this Hessian identity gives
\begin{align*}
	\Delta X = - \mathrm{Ric} ^{\sharp}(X).
\end{align*} 
The Killing equation gives $(d \alpha )_{ij} = 2 \nabla _{i} \alpha _j$, and thus
\begin{align*}
	(\delta d \alpha )_{j} = - \nabla ^{i}(d \alpha )_{ij} = 2 \mathrm{Ric}_{jk} X ^{k},
\end{align*} 
where $\delta = - * d *$ is the codifferential. Applying Hodge star gives
\begin{align*}
	d (* d \alpha ) = 2 * \mathrm{Ric}_g(X, \cdot ).
\end{align*} 
Finally, use the facts that $\omega _{X} = \iota _{X}(* d \alpha )$, $\mathcal{L}_{X} (* d \alpha ) =0$ and $\mathcal{L}_{X} = d \iota _{X} + \iota _{X} d$.
\end{rmk}
\begin{comment}
\begin{rmk}
	By Frobenius' theorem, wherever $X\ne0$, $\omega_X=0$ if and only if there exist local hypersurfaces perpendicular to $X$. When $X \neq 0$, consider the three-dimensional distribution $\mathcal{H}_{X} = X ^{\perp} = \mathrm{ker} \alpha $. If $V, W$ are vector fields perpendicular to $X$, then $d \alpha (V, W) = - \alpha ([V,W])$. Thus, $d \alpha $, restricted to $ \mathcal{H}_{X}$, measures how moving in two perpendicular-to-$X$ directions can produce a displacement in the $X$-direction. With our conventions, whenever $X \neq 0$, $d \omega _X =0 $ if and only if $\mathrm{Ric}_g(X, V) =0$ for all $V \in \mathcal{H}_{X}$.
\end{rmk}
\end{comment}

By the Einstein assumption, $\omega _{X}, \omega _{Y} $ are closed, i.e.
\begin{align*}
	d \omega _X =0, \ \ d \omega _{Y} =0.
\end{align*} 
Since $T^2$ is abelian, $X, Y$ commute, and consequently $\mathcal{L}_{Y} \omega _X = \mathcal{L}_{X} \omega _{Y} =0$. By Cartan's formula, $d (\iota _{Y} \omega _X) = d (\iota _{X} \omega _Y)=0$, thus the scalars $\iota _{Y} \omega _X$ and $\iota _{X} \omega _Y$ are constants. Since $X$ and $Y$ vanish at some point, we know $\iota _{Y} \omega _X = \iota _{X} \omega _Y \equiv 0$. Since $\iota _{X} ( * \eta ) = * ( \eta \wedge \alpha )$ for a differential form $\eta$, we know
\begin{align*}
	\alpha \wedge \beta \wedge d \alpha =0,\ \ \alpha \wedge \beta \wedge d \beta =0.
\end{align*} 
Consider the distribution $\mathcal{H} = \mathrm{ker} \alpha \cap \mathrm{ker} \beta = ( \mathrm{span} \{X, Y\} ) ^{\perp}$ on the principal region $M^{\circ}:=\pi^{-1}(B^{\circ})$. Then $d \alpha = d \beta =0 $ on $\Lambda ^2 \mathcal{H}$. Frobenius' theorem shows that $\mathcal{H}$ is integrable. 

We now express this in connection notation. Let $(X_1,X_2)=(X,Y)$ and $G_{ij}=g(X_i,X_j)$. In the original angular coordinates $\theta_0^i=\arg z_i$, every $T^2$-invariant metric on $M^\circ$ can be written as
\begin{align*}
g &= h+G_{ij}(d\theta_0^i+A^i)(d\theta_0^j+A^j),\\
A^i &= A^i_\mu\,dx^\mu ,
\end{align*}
where $x^\mu$ are local coordinates on $B^\circ$ and the one-forms $A^i$ on the base are pulled back to $M^\circ$. The orthogonal torus connection has components
\[
\eta^i:=d\theta_0^i+A^i=G^{ij}g(X_j,\cdot),
\qquad \eta^i(X_j)=\delta^i_j.
\]
They form the dual coframe to $(X_1,X_2)$ on the vertical distribution, extended to vanish on $\mathcal H$. The metric dual forms satisfy
\[
\alpha=G_{11}\eta^1+G_{12}\eta^2,\qquad
\beta=G_{12}\eta^1+G_{22}\eta^2.
\]
Writing $A=A^i e_i$ in a basis of the Lie algebra, the general curvature formula is $F=dA+\tfrac12[A\wedge A]$. The Lie algebra of $T^2$ has zero bracket, so $F^i=dA^i$. The horizontal lifts satisfy
\[
U_\mu=\partial_{x^\mu}-A^i_\mu X_i,
\qquad
[U_\mu,U_\nu]=-F^i_{\mu\nu}X_i.
\]
Thus the integrability proved above gives $F^i=0$.

The original circle-valued angles trivialize the principal region, so the $A^i$ are globally defined on $B^\circ$. Since $B^\circ$ is simply connected, $dA^i=0$ implies $A^i=df^i$ for smooth functions $f^i$ on $B^\circ$. Define adapted angles by $\theta^i=\theta_0^i+f^i$. Then $d\theta_0^i+A^i=d\theta^i$, and the new connection potentials are $A^i-df^i=0$. Consequently,
\[
g=h+G_{ij}\,d\theta^i d\theta^j,
\qquad
\partial_{\theta^i}=X_i.
\]
Here the adapted angles $\theta^i$ take values in $\mathbb R/2\pi\mathbb Z$.

Note that $G_{ij}$ is a globally defined matrix-valued function on $B ^{\circ}$ : for any $q \in B ^{\circ}$ and any $p \in \pi ^{-1}(q)$, $G_{ij}(q) := g_{p}(X_i(p), X_j(p) )$ is independent of the choice of $p$. Set
\begin{align*}
	\rho := \sqrt{\mathrm{det} G} .
\end{align*} 
Let $K$ be the Gaussian curvature of $h$. Let $\mathcal{V} = \mathcal{H}^{\perp}$ be the vertical distribution.

By the Koszul formula, for any basic horizontal fields $U, V$, the connection formulas are
\begin{align*}
	\nabla ^{g}_{U} V &= \nabla ^{h}_{U} V,\\
	\nabla ^{g}_{U} X_i = \nabla ^{g}_{X_i} U &= \frac{1}{2} G ^{j k} U(G_{k i}) X_j, \\
	\nabla ^{g}_{X_i} X_j &= -\frac{1}{2} \nabla ^{h} G_{ij}.
\end{align*} 
Let $\mu,\nu$ denote horizontal coordinate indices, write $G_\mu=\partial_\mu G$, and let $G_{ij;\mu\nu}$ denote the $h$-Hessian of the scalar component $G_{ij}$. The curvature tensor relation is
\begin{align*}
R_{i j k l} &= \frac{1}{4} \left( \langle d G_{il}, d G_{jk} \rangle_{h} - \langle d G_{ik}, d G_{jl} \rangle_{h} \right), \\
	R_{\mu i \nu j} &= - \frac{1}{2} G_{ij; \mu \nu } + \frac{1}{4} (G_{\nu } G ^{-1} G_{\mu }) _{ij},\\
	R_{\mu \nu i j} &= \frac{1}{4}\left( G_{\nu } G ^{-1} G_{\mu } - G_{\mu } G^{-1} G_{\nu } \right) _{ij}, \\
	R_{\mu \nu \kappa \lambda }&= K (h_{\mu \kappa} h_{\nu  \lambda } - h_{\mu \lambda } h_{\nu \kappa}),\\
	R_{\mu \nu \kappa i} &= R_{\mu i j k} =0,
\end{align*} 
the Ricci curvature relation is
\begin{align*}
	\mathrm{Ric}_{ij} &= -\frac{1}{2} \Delta _{h} G_{ij} + \frac{1}{2} h ^{\mu \nu }\left( G_{\mu } G ^{-1} G_{\nu } \right) _{ij} - \frac{1}{2} \langle d \log \rho , d G_{ij} \rangle_{h},\\
	\mathrm{Ric}_{\mu \nu } &= K h_{\mu \nu } - \nabla ^2_{\mu \nu }(\log \rho )- \frac{1}{4} \mathrm{tr}(G ^{-1}G_{\mu } G ^{-1} G_{\nu }),\\
	\mathrm{Ric}_{\mu i} &=0,
\end{align*} 
and the scalar curvature relation is
\begin{align*}
	R_g &= 2K - 2\Delta _h (\log \rho ) - |d \log \rho |^2_h - \frac{1}{4} h ^{\mu \nu } \mathrm{tr}(G ^{-1} G_{\mu } G ^{-1} G_{\nu }).
\end{align*} 

The vertical part of Ricci curvature is equivalent to
\begin{align*}
	\mathrm{Ric} ^{\mathcal{V}} = - \frac{1}{2 \rho } G \mathrm{div}_{h} \left( \rho G ^{-1} d G \right) ,
\end{align*} 
whose trace gives
\begin{align*}
	\mathrm{tr}(G ^{-1} \mathrm{Ric} ^{\mathcal{V}}) = - \frac{\Delta _h \rho }{\rho }.
\end{align*} 

Using the Einstein equation $\mathrm{Ric} ^{\mathcal{V} } = 3 G$, for $H:= \rho ^{-1}G$ with $\mathrm{det} H =1$, we obtain
\begin{align*}
	\Delta _{h} \rho = -6 \rho ,\ \ \mathrm{div}_{h}(\rho H ^{-1} d H ) =0.
\end{align*} 
Expanding this matrix equation, we obtain
\begin{align*}
	\mathrm{div}_{h}(\rho d H_{12}) - \rho e H_{12} =0,
\end{align*} 
where
\begin{align*}
	e:= |d H_{12}|^2_{h} - \langle d H_{11}, d H_{22} \rangle_{h} = |d H_{12} - H_{12} d \log H_{11}|^2_{h} + |d \log H_{11}|^2_{h} \geq 0.
\end{align*} 
We now use this equation of $H_{12}$ together with its boundary values to conclude that $H_{12} \equiv 0$. Near the boundary arc corresponding to $z_i=0$, consider small neighborhoods as the orbit space of $r_i=|z_i| \ll 1$. Then $X_i = r_i E_i$ with $0<c \leq |E_i|_g \leq C $ for some positive constants $c,C>0$, which can be shown by the fact that $|E_i|_{g_0}=1$ for the round metric and the equivalence between $g$ and $g_0$. The fields $E_1,E_2$ are orthonormal for the round comparison metric, so their entire $g$-Gram matrix is bounded below by $cI$, and its determinant is at least $c^2$. At an ordinary axis the collapsing-circle half-turn reverses its normal plane and fixes the tangent plane, forcing the limiting cross term to vanish; at a corner the two distinct torus rotation planes are orthogonal. Smoothness then gives $g(E_1,E_2)\to0$ along every approach to the boundary. Thus
\begin{align*}
	H_{12} = \frac{g(E_1, E_2)}{\sqrt{g(E_1, E_1) g(E_2, E_2) - g(E_1, E_2) ^2} } \to 0 \text{ as } r_i\to 0.
\end{align*} 
Thus, by defining $H_{12}$ as zero on the boundary $\partial B$, we obtain
\begin{align*}
	H_{12} \in C^{\infty}(B ^{\circ}) \cap C ^{0}(B) ,\ \ H_{12}|_{\partial B} =0.
\end{align*} 
By the strong maximum principle, we have $H_{12} \equiv 0$. Thus, $G_{12} \equiv 0$ in $B ^{\circ}$, and, writing $(\phi,\theta)=(\theta^1,\theta^2)$ for the adapted angles above, we can write
\begin{align*}
	g = h + a ^2 (d \phi )^2 + b ^2 (d \theta )^2,
\end{align*} 
for some smooth functions $a, b>0$ on $B ^{\circ}$.

\subsection{Curvature upper bound in the regular part}
For the diagonal metric obtained above, the curvature formulas simplify. On the principal region $M^\circ$,
\begin{align*}
	g = h + a ^2 (d \phi )^2 + b ^2 (d \theta )^2 \text{ on } M ^{\circ},
\end{align*}
with $a, b >0$ on $B ^{\circ}$ and $b = 0 $ on  $\Gamma _1$ and $a = 0$ on $\Gamma _2$.
Recall $\rho = a b$.
Let $E_{\phi }, E_{\theta }$ be vertical orthonormal fields, and $U$ be any unit horizontal field. Then
\begin{align*}
	\mathrm{sec}_g|_{\mathcal{H}} &= K,\\
	\mathrm{sec}_g(U, E_{\phi }) &= - \frac{(\nabla ^2_h a) (U, U)}{a},\\
	\mathrm{sec}_g(U, E_{\theta }) &= - \frac{(\nabla ^2_h b)(U, U)}{b},\\
	\mathrm{sec}_g(E_{\phi }, E_{\theta }) &= - \frac{\langle d a, d b \rangle_{h}}{ab},
\end{align*} 
and
\begin{align*}
	\mathrm{Ric}_g|_{\mathcal{H}} &= K h - \frac{\nabla ^2_h a}{a} - \frac{\nabla ^2_h b}{b}, \\
	\mathrm{Ric}_g(E_{\phi }, E_{\phi }) &= - \frac{\Delta _h a}{a} - \frac{\langle d a , d b \rangle_h}{a b},\\
\mathrm{Ric}_g(E_{\theta }, E_{\theta }) &= - \frac{\Delta _h b}{b} - \frac{\langle d a , d b \rangle_h}{a b},\\
\mathrm{Ric}_g(E_{\phi }, E_{\theta }) &=\mathrm{Ric}_g(E_{\phi }, U) = \mathrm{Ric}_{g}(E_{\theta }, U)=0,
\end{align*} 
and
\begin{align*}
	R_g &= 2K - \frac{2 \Delta _h a}{a} - \frac{2 \Delta _h b}{b} - \frac{2 \langle d a, d b \rangle_h}{ab}.
\end{align*} 
Within this subsection, unmarked derivatives and tensor norms are taken with respect to $h$.

Using $\mathrm{Ric}_g = 3 g$, we have
\begin{align*}
	\frac{\nabla ^2 a}{a} + \frac{\nabla ^2 b}{b} &= (K-3) h,\\
	\frac{\Delta a}{a} + \frac{\langle d a, d b \rangle_h}{a b} &= \frac{\Delta b}{b} + \frac{\langle d a, d b \rangle_h}{a b} = -3.
\end{align*} 
So 
\begin{align*}
	K = - \frac{\langle d a, d b \rangle_h}{a b},\ \ \Delta \rho = - 6 \rho .
\end{align*} 
Applying the Bochner formula to $(B ^{\circ}, h)$,
\begin{align*}
	\Delta  \langle d a, d b \rangle &= \langle d \Delta a, d b \rangle + \langle d a, d \Delta b \rangle + 2 \langle \nabla ^2 a, \nabla ^2 b \rangle + 2K \langle d a, d b \rangle.
\end{align*} 
Therefore,
\begin{align*}
	\Delta K + 3 \langle d \log \rho , d K \rangle= 3 (K-1) (K+3)+ \frac{1}{2} \left| a ^{-1} \nabla ^2 a - b ^{-1} \nabla ^2 b \right|^2.
\end{align*} 

We prove $K\le1$ on $B^\circ$ by a cutoff argument.

First, by the curvature relation, we know that $|K| \leq \max| \mathrm{sec}_g| \leq  C(g)$ in $B ^{\circ}$ for some constant $C(g)$ depending on $g$. Second, $\rho = | X \wedge Y|_{g} >0 $ is well defined on $M ^{\circ}$ and can be extended as a global Lipschitz function on $M$ so that $|\nabla_g \rho |\leq C(g)$ and $\rho =0$ precisely on $M \setminus M ^{\circ}$.

Note that for any $T^2$-invariant function $f$ on $M ^{\circ}$,
\begin{align*}
	\Delta _g f = \Delta _h f + \langle d \log \rho , d f \rangle_h.
\end{align*} 
We can rewrite the equation of $K$ as
\begin{align*}
	\mathrm{div}_g( \rho ^2 \nabla _g K) &= \rho ^2 \left(  3 (K-1) (K+3)+ \frac{1}{2} \left| a ^{-1} \nabla ^2_h a - b ^{-1} \nabla ^2_h b \right|_h^2\right) \\
					     &\geq 3 \rho ^2 (K-1)(K+3).
\end{align*} 
Choose a cutoff function $\eta $ so that $\eta = 0$ on $[0,1]$, $\eta =1$ on $[2, \infty)$ and $0\leq \eta ' \leq 2$. For $0<\epsilon \ll 1$, define
\begin{align*}
	\chi _{\epsilon } := \eta \left( \frac{\rho }{\epsilon } \right) .
\end{align*} 
Then $\mathrm{supp} \chi _{\epsilon } \subset M ^{\circ}$.
Set $w= (K-1)_{+}$. Multiplying this inequality by $\chi _{\epsilon }^2 w$, and integrating by parts, we obtain
\begin{align*}
	3 \int_{M ^{\circ}} \chi _{\epsilon }^2 \rho ^2 w ^2 (K+3) &\leq - \int_{M ^{\circ}} \langle \rho ^2 \nabla _g K, \nabla_g (\chi _{\epsilon }^2 w) \rangle_g \\
	&= - \int_{M ^{\circ}} \rho ^2 \chi ^2_{\epsilon } |\nabla _g w|^2 - 2 \int_{M ^{\circ}} \rho ^2 \chi _{\epsilon } w \langle \nabla _g w, \nabla _g \chi _{\epsilon } \rangle,
\end{align*} 
which gives
\begin{align*}
		\frac{1}{2}\int_{M ^{\circ}} \rho ^2 \chi ^2_{\epsilon } |\nabla _g w|^2 + 3 \int_{M ^{\circ}} \chi _{\epsilon }^2 \rho ^2 w ^2 (K+3) &\leq 2 \int_{M ^{\circ}} \rho ^2 w ^2 |\nabla_g \chi _{\epsilon }|^2 \\
																		      &\leq C(g) \int_{\{\epsilon \leq \rho \leq 2 \epsilon \} } w ^2 \\
																		      &\leq C(g) \mathrm{Vol}( \{\epsilon \leq \rho \leq 2 \epsilon \} ) \to 0,
\end{align*} 
as $\epsilon \to 0$. On the other hand, $\chi _{\epsilon } \to 1$ as $\epsilon \to 0$, so we obtain
\begin{align*}
	\int_{M ^{\circ} \cap \{K> 1\} } \rho ^2 w ^2 (K+3) =0,
\end{align*} 
which proves the claim that $K \leq 1$ in $B ^{\circ}$.

\subsection{Convex exhaustion criterion}

\begin{prop}\label{prop-convex-exhaustion}
	Let $(\Omega , h)$ be a connected smooth Riemannian manifold without boundary, not necessarily complete. Suppose $\mathrm{sec}_{h} \leq 1.$ 
Assume that there exists $F \in C ^{\infty}(\Omega )$ such that:
\begin{align*}
	\{F \leq r\} \text{ is compact for every } r \in \mathbb{R},\ \ \nabla ^2_{h} F >0.
\end{align*} 
Let $d_h$ denote the induced length distance on $(\Omega, h) $, that is, $d_{h}(x,y) = \inf \{L_h(\gamma ): \gamma \subset \Omega \text{ is a smooth curve between } x \text{ and } y\} $. Then both $(\Omega , d_{h})$ and its metric completion are $\mathrm{CAT}(1)$.
\end{prop}

\begin{proof}
	Fix $p_0 \in \Omega $. Choose regular values $r_j \to \infty$ so that $p_0 \in \{F < r_1\} $. Let $C_j$ be the connected component of $\{F \leq r_j\} $ containing $p_0$, and let $d_j$ be the induced length metric on $C_j$. Then $C_j \subset C_{j+1}$ and we have an exhaustion $\Omega = \cup _{j} C_j$.

	Each $C_j$ is locally geodesically convex. Indeed, for any $p \in C_j$, we can find a sufficiently small strongly convex neighborhood $V \subset \Omega $. By the curvature assumption, we can assume $(V, d_h)$ is $\mathrm{CAT}(1)$. For any $x, y \in V \cap C_j$, let $\gamma $ be the constant-speed minimizing $h$-geodesic. Since $\nabla ^2_h F >0$, we know $F(\gamma (t) ) \leq (1-t) F(x) + t F(y) \leq r_j$. Therefore, $V \cap C_j$ is geodesically convex in $V$, and $d_j(x, y) = d_h(x, y)$ for all such $x, y \in V \cap C_j$. Thus $(C_j,d_j)$ is a compact, locally $\mathrm{CAT}(1)$ length space. Each of its local geodesics is locally an ambient $h$-geodesic, by the local metric agreement. Along a nonconstant closed such geodesic $\gamma$, strict convexity would give $(F\circ\gamma)^{\prime\prime}>0$, contradicting periodicity. Hence there are no closed local geodesics. The compact globalization criterion makes $(C_j,d_j)$ globally $\mathrm{CAT}(1)$ \cite[Corollary~3.3.4]{notes-AKP}. 

	For $j'\ge j$, a minimizing $d_{j'}$-geodesic between points of $C_j$ is locally an ambient geodesic. Convexity of $F$ keeps it in $\{F\le r_j\}$, hence in the same component $C_j$. Thus $d_{j'}=d_j$ on $C_j$. Every ambient comparison curve lies in some $C_k$, so $d_h=d_j$ on $C_j$ as well. For any geodesic triangle in $(\Omega , d_h)$, one can choose large enough $j$ so that it lies in $C_j$, and thus $(\Omega , d_h)$ is $\mathrm{CAT}(1)$. Finally, by Bridson--Haefliger \cite[Corollary II.3.11]{notes-BH}, metric completion preserves $\mathrm{CAT}(1)$.
\end{proof}

In our case, $\{ - \log \rho\leq  r\} \Subset B ^{\circ}$ is compact and gives an exhaustion of $B ^{\circ}$. Moreover,
\begin{align*}
	\nabla ^2_h (- \log \rho)   &= (3-K) h + d \log a \otimes d \log a + d \log b \otimes d \log b \\
		      & \geq 2 h>0.
\end{align*} 
Consequently, both $(B ^{\circ}, d_h)$ and its completion $\widehat{(B ^{\circ}, d_h)}$ are $\mathrm{CAT}(1)$. 

Let $(B, d_B)$ denote the orbit distance, that is,
\[
d_B( \pi (p), \pi (q) ) = \min_{\tau \in T^2}d_g(p, \tau (q) ).
\]
We claim that $(B,d_B)=\widehat{(B^\circ,d_h)}$, so Proposition~\ref{prop-convex-exhaustion} makes $(B,d_B)$ a $\mathrm{CAT}(1)$ space. It suffices to show that $(B ^{\circ}, d_h) \subset (B, d_B)$ is an isometric inclusion. 

Projection of curves gives $d_B\le d_h$. To prove the reverse inequality, take two points of $B^\circ$ and choose representatives $p,q\in M^\circ$ so that $d_{B}(\pi (p), \pi (q) ) = d_{g}(p,q)$ and let $\gamma \subset (M, g)$ be a minimizing geodesic between $p$ and $q$. Then $\gamma (s) = (z_1(s), z_2(s), t(s) ) $ may have intersection with singular orbits $\{ (z_1, z_2, t): z_1 z_2 =0\} $. We can perform a slight perturbation so that each $z_i(s) \in \mathbb{C}$ avoids $0 \in \mathbb{C}$ and therefore obtain perturbations $\gamma _n \subset M ^{\circ}$ with the same endpoints so that $\gamma _n \to \gamma $ in $C^1$. The conclusion then follows from $d_h(\pi (p), \pi (q) ) \leq L_{h}(\pi (\gamma _n) ) \leq L_{g}(\gamma _n) \to d_{g}(p, q) $.

\subsection{Maximal diameter and rigidity}
\begin{lem}\label{lem-fixed-quotient}
	Let a compact group $G$ act isometrically on a closed Riemannian manifold $(M, g)$. Let $F$ be a connected component of $\mathrm{Fix}(H)$, for some subgroup $H \subset G$, and assume $F$ is $G$-invariant.
Give $F$ its induced Riemannian metric and $F / G$ its intrinsic quotient distance. Then
$F / G \subset M / G$ 
locally preserves distances.
\end{lem}
\begin{proof}
	Take two orbit points $x, y \in F / G$ whose distance in $M / G$ is less than $\mathrm{inj}(M,g)$, and take $p, q \in F$ so that $d_{M / G}(x, y) = d_g(p, q)$. Let $\gamma \subset (M,g)$ be the unique minimizing geodesic between $p$ and $q$. So for any $h \in H$, $h\circ \gamma = \gamma $, thus $\gamma \subset F$ and $d_{F}(p,q) = d_g(p, q)$. Therefore, $d_{F / G}(x, y) \leq d_{M / G}(x, y)$. The opposite inequality always holds.
\end{proof}

In our case, let $F_1= \{z_2=0\} $ and $F_2= \{ z_1=0\} $ so that $\pi (F_i) = \Gamma _i$, $H_1 = \{1\} \times S^1$ and $H_2= S^1 \times \{1\}$. Then $F_i=\mathrm{Fix}(H_i)$, so Lemma~\ref{lem-fixed-quotient} makes each $\Gamma_i=F_i/T^2\simeq[0,\ell_i]\subset B$ a local geodesic. 

By Bridson--Haefliger \cite[Proposition II.1.4]{notes-BH}, inside a global $\mathrm{CAT}(1)$ space, every local geodesic of length less than $\pi $ is minimizing, and two points at distance less than $\pi $ have a unique minimizing geodesic. Therefore, $\max \{ \ell_1, \ell_2 \} \geq \pi $. An arc of length at least $\pi$ has initial subarcs of every length $s<\pi$, each minimizing by the same CAT property. Letting $s\uparrow\pi$ gives $\mathrm{diam}(B)\ge\pi$.

Since $\pi $ is $1$-Lipschitz, together with Bonnet-Myers, we have 
$$\pi \leq \mathrm{diam}(B) \leq \mathrm{diam}(M, g) \leq \pi .$$
Thus, $\mathrm{diam}(M,g)=\pi$. Cheng's maximal-diameter rigidity \cite{notes-Cheng} proves Theorem~\ref{thm-sphere}.

\section{\texorpdfstring{Complete Ricci-flat metrics on $\mathbb R^4$}{Complete Ricci-flat metrics on R4}}\label{sec-R4}
This section proves Theorem~\ref{notes-R4-classification}.

\subsection{Model spaces}
The two models in Theorem~\ref{notes-R4-classification} admit the following explicit expressions. On $\mathbb C^2$, write
\begin{align*}
	\zeta _1 = x e ^{i \theta ^1},\ \  \zeta _2 = y e ^{i \theta ^2},\ \  x, y \geq 0,\ \ \theta ^{j} \in \mathbb{R} / (2 \pi \mathbb{Z}).
\end{align*} 
Green functions in these models are functions on the four-manifold, normalized by $-\Delta_gG_p=4\pi^2\delta_p$. The torus translates the two angles, with generators $X_j = \partial _{\theta ^{j}}$. Its quotient is the quadrant with coordinates $\{(x, y): x, y \geq 0 \}$. Set $r^2=x^2+y^2$. The Hodge-star computations below are on the smooth quotient interior $x,y>0$; the displayed scalar formulas extend continuously to the axes.

\subsubsection{Euclidean space}
The metric, quotient metric and orbit Gram matrix are
\begin{align*}
	g_0 &= d x ^2 + d y ^2 + x ^2 (d \theta ^1) ^2 + y ^2  (d \theta ^2) ^2,\\
	h_0 &= d x ^2 + d y ^2,\ \ G_0 = \begin{pmatrix} x ^2 & 0 \\ 0 & y ^2 \end{pmatrix}. 
\end{align*} 
Thus $$\rho = \sqrt{\mathrm{det} G_0} = x y .$$
Since $* _{h_{0}} d x = d y, *_{h_0} d y = - d x$,
 \begin{align*}
	d \tau = - *_{h_0} d (x y) = x d x - y d y,\ \ \tau = \frac{x ^2 - y ^2}{2}.
\end{align*} 
For
\begin{align*}
	\xi := \sqrt{\rho ^2 + \tau ^2} = \frac{x ^2+ y ^2}{2},
\end{align*} 
we have
\begin{align*}
	x = \sqrt{\xi + \tau } ,\ y= \sqrt{\xi - \tau } .
\end{align*} 
Thus $(\rho,\tau)$ identifies the closed quotient with the half-plane $\rho\ge0$, conformally on its interior; equivalently $\tau + i \rho = \frac{1}{2}(x+ i y)^2$, and
\begin{align*}
	d \rho ^2 + d \tau ^2 = 2 \xi h_0.
\end{align*} 

The normalized Green function of $g_0$ at the origin  is
\begin{align*}
	G_p = \frac{1}{2 \xi } = \frac{1}{r^2},
\end{align*} 
and the full Weyl expression is
\begin{align*}
	g_0 = \frac{d \rho ^2+ d \tau ^2}{2\xi } + (\xi +\tau ) (d \theta ^{1})^2 + (\xi - \tau ) (d \theta ^2)^2.
\end{align*} 

\subsubsection{Taub--NUT}
Set
\begin{align*}
	\eta = x ^2 d \theta ^{1} + y ^2 d \theta ^2,\ \ f(r)= 1+ \frac{1}{4}r ^2,\ \ b(r) = \frac{1}{4}+ \frac{1}{4 f(r)}.
\end{align*} 
Then the Taub--NUT metric can be written as
\begin{align*}
	g_{\mathrm{TN}} = f g_0 - b \eta ^2.
\end{align*} 
The metric has eigenvalues $f, f, f, f ^{-1}$ relative to $g_0$. The quotient metric and orbit Gram matrix are
\begin{align*}
	h_{\mathrm{TN}} = f (d x ^2 + d y ^2) ,\ \ G_{\mathrm{TN}} = \begin{pmatrix} f x ^2- b x ^{4} & -b x ^2 y ^2 \\ - b x ^2 y ^2 & f y ^2 - b y ^{4} \end{pmatrix} .
\end{align*} 
Since $f - b r ^2 = f ^{-1}$, 
\begin{align*}
	\rho _{\mathrm{TN}} = \sqrt{\mathrm{det} G_{\mathrm{TN}}} = xy.
\end{align*} 
Since $h_{\mathrm{TN}}=fh_0$ and the Hodge star on one-forms is conformally invariant in dimension two,
\begin{align*}
	d \tau = - * _{h_{\mathrm{TN}}} d (x y) = x d x - y d y,\ \ \tau = \frac{x ^2 - y ^2}{2},\ \ \xi = \frac{x ^2+ y ^2}{2}.
\end{align*} 
The same half-plane coordinates therefore give
\begin{align*}
	h_{\mathrm{TN}} = \left( \frac{1}{2 \xi }+ \frac{1}{4} \right) (d \rho ^2+ d \tau ^2).
\end{align*} 
Here $r$ is the coordinate radius on $\mathbb C^2$. The normalized Green function on $(\mathbb R^4,g_{\mathrm{TN}})$ is again $r^{-2}$.

\subsection{Action reduction}
Let $M^4$ be a smooth $4$-manifold homeomorphic to $\mathbb{R}^{4}$, which admits an effective smooth $T^2$-action.
By Orlik--Raymond, the $T^2$-symmetry forces its smooth structure to be standard, and up to a diffeomorphism of $\mathbb{R}^{4}$ and an automorphism of $T^2$-action, we can assume that the effective $T^2$-action
is given by the standard action:
\begin{align*}
	\left( e ^{i s_1}, e ^{i s_2}\right) \cdot (z_1, z_2) = (e ^{i s_1} z_1, e ^{i s_2} z_2).
\end{align*} 

To see this, extend the action continuously to $\hat M=M\cup\{\infty\}\simeq S^4$ by fixing $\infty$. Apply the continuous-action appendix of Orlik--Raymond with $A=\varnothing$, since $\hat M\simeq S^4$ is an orientable integral cohomology four-manifold. Since $H_1(\hat M;\mathbb Z)=0$, their classification gives a disk quotient with no exceptional finite-isotropy orbits. The fixed-point count is $\chi(\hat M)=2$. Its two boundary arcs have primitive isotropy generators forming an integral basis. Deleting $\infty$ leaves a quadrant with one corner and two rays.

This compactification is used only for the topological orbit data. On the original smooth $M$, the smooth slice models make the quotient a smooth surface with corners. Relative uniqueness of smooth structures on surfaces identifies it with the standard quadrant, preserving the labels and corner chart \cite[Theorem~B and relative-boundary version]{notes-Hatcher}; smooth action reconstruction then gives the asserted equivariant diffeomorphism. This reconstruction applies also to noncompact manifolds \cite[\S\S1.11--1.12, p.~537, and Appendix]{notes-OrlikRaymond}. 

Now let $g$ be a smooth complete Ricci-flat metric preserved by the action. Let $\pi:M\to X$ be the quotient map and $p\in M$ the unique torus fixed point, projecting to the corner. Write $X^\circ\simeq(0,\infty)^2$ and $M^\circ=\pi^{-1}(X^\circ)$. Let $X_1,X_2$ be the primitive $2\pi$-periodic Killing generators and $\alpha_i$ their metric dual forms.

Since $X ^{\circ}$ is contractible, the free orbit bundle is trivial, and in standard angular coordinates $\theta _0 ^{i}$, a general $T^2$-invariant metric has the form
\begin{align*}
	g = h + G_{ij} (d \theta ^{i}_{0} + A ^{i})(d \theta _0 ^{j} + A ^{j}).
\end{align*} 

The twist argument in Section~\ref{subsec-orthogonal} applies because both generators vanish at $p$. Thus $\mathcal{H}=\ker\alpha_1\cap\ker\alpha_2$ is integrable. So $d A ^{i} =0$, and contractibility gives $A ^{i} = d f ^{i}$ globally, and the change $\theta ^{i} = \theta _0 ^{i} + f ^{i}$ modulo $2 \pi $ gives
\begin{align*}
	g = h + G_{ij} d \theta ^{i} d \theta ^{j} \text{ on } M ^{\circ}.
\end{align*} 
Let $\rho = \sqrt{\mathrm{det}G} $ be the fiber density function. Then $\rho >0$ in $X ^{\circ}$ and $\rho =0$ on $\partial X$. 

\subsection{Global Weyl coordinates}
Using the Ricci-flat equation $\mathrm{Ric}_g = 0$, we obtain
\begin{align*}
	\Delta _h \rho =0.
\end{align*} 

We can construct the harmonic conjugate as follows.
Fix an orientation of the smooth surface $(X^\circ,h)$. The one-form $-*_h d\rho$ is closed, since $d(-*_h d\rho)=-(\Delta_h\rho)dV_h=0$. As $X^\circ$ is simply connected, there is a smooth function $\tau:X^\circ\to\mathbb R$, unique up to an additive constant, with $d\tau=-*_h d\rho$. Thus, 
$$\Delta _{h} \tau =0,\ \ \langle d \rho , d \tau  \rangle_{h} =0,\ \ |d \tau |_{h} = |d \rho |_{h}.$$ 

The goal is to show that a continuous extension of $(\rho , \tau )$ gives a global coordinate of $X$.

\begin{lem}\label{lem-conjugate}
	The lift $\tilde{\tau }=\tau \circ \pi $ extends smoothly and harmonically to $M$, and can be normalized by $\tilde{\tau} (p)=0$. Also, $\tilde{\tau }$ descends to a continuous function on $X$, which is still denoted by $\tau $. Define $\xi := \sqrt{\rho ^2+ \tau ^2} $. Near $p$, we have the expansion
	\begin{align*}
		\xi = \frac{1}{2}d ^2+ O(d ^{4}),
	\end{align*} 
	where $d = d_g(p, \cdot )$.
\end{lem}
\begin{proof}
	On $M ^{\circ}$, we have
	\begin{align*}
		\Delta _{g} \tilde{\tau } = ( \Delta _{h} \tau + \langle d \log \rho , d \tau  \rangle_h ) \circ \pi =0,\ \ |\nabla \tilde{\tau }|_g = |\nabla \tilde{\rho }|_g,
	\end{align*} 
	where $\tilde{\rho } = \rho \circ \pi = |X_1 \wedge X_2|_g$. So $|\nabla \tilde{\tau }|_g$ is locally bounded.

	The singular set is the union of two smooth codimension-two fixed surfaces meeting at $p$. In a sufficiently small geodesically convex ball $U$, a geodesic between principal points can be approximated, with fixed endpoints and converging lengths, by smooth curves avoiding both surfaces. Indeed, a transverse curve misses each codimension-two stratum and the isolated point. Integrating the locally bounded gradient along these curves gives
\[
 |\tilde\tau(x)-\tilde\tau(y)|\le L_Ud_g(x,y),\qquad x,y\in U\cap M^\circ.
\]
Thus $\tilde\tau$ extends uniquely as a locally Lipschitz function on $M$, and belongs to $W^{1,2}_{\mathrm{loc}}(M)$. Continuity preserves torus invariance, so the extension descends continuously to $X$.

	As in the $S^4$-case, let $\chi _{\epsilon } = \eta (\epsilon ^{-1} \rho )$ be a cutoff function for some $\eta $ so that $\eta =0$ on $[0,1]$, $\eta =1$ on $[2, \infty)$ and $0 \leq \eta ' \leq 2$. For any test function $\varphi \in C ^{\infty}_{c}(M)$, we have
\begin{align*}
	\int_M\chi_\epsilon\langle\nabla^g\tilde\tau,\nabla^g\varphi\rangle_g\,dV_g
 &= -\int_M\varphi\langle\nabla^g\tilde\tau,\nabla^g\chi_\epsilon\rangle_g\,dV_g=0.
\end{align*} 
Here the cutoff uses the lift of $\rho$. Orthogonality of $d\tau$ and $d\rho$ makes the last integral zero. The singular set has zero Riemannian volume, so dominated convergence proves weak harmonicity on $M$. Elliptic regularity makes $\tilde\tau$ smooth.

Normalize $\tilde\tau(p)=0$. The action reduction has fixed a primitive torus basis whose tangent representation at $p$ consists of two standard rotation planes. These planes are orthogonal for the invariant inner product $g_p$. Choose orthonormal coordinates on them and use the exponential map, which intertwines this linear action with the given action. The resulting local geodesic normal coordinates $(u_1,u_2,u_3,u_4)$ satisfy $g_{\alpha\beta}=\delta_{\alpha\beta}+O(|u|^2)$ and the exact formulas
\begin{align*}
	X_1 = - u_2 \partial _{u_1} + u_1 \partial _{u_2},\ \ X_2= - u_4 \partial _{u_3}+ u_3 \partial _{u_4}.
\end{align*} 
Set
\begin{align*}
	x_p^2=u_1^2+u_2^2,\qquad y_p^2=u_3^2+u_4^2,\qquad d^2=d_g(p,\cdot)^2=|u|^2=x_p^2+y_p^2.
\end{align*} 
In this local computation, abbreviate $x=x_p$ and $y=y_p$; global action radii will be specified separately.

Since $G_{ij}$ is $T^2$-invariant, we can focus on one slice $(u_1, u_2, u_3, u_4) = (x, 0, y, 0)$. Then 
\begin{align*}
	G_{11} &= \langle X_1, X_1 \rangle_g = x ^2 g_{22} = x ^2(1+ O(d ^2) ),\\
	G_{22} &= \langle X_2, X_2 \rangle_g = y ^2 g_{44} = y ^2(1+ O(d ^2) ),\\
	G_{12} &= x y g_{2 4} = x y \cdot  O(d ^2).
\end{align*} 
In fact, for later use, we have better asymptotics of $G_{12}$. Note that $g_{24}(0, 0, y, 0) = g_{24}(x, 0, 0, 0)=0$. To see this, one can consider the isometric rotation by angle $\pi $ in the plane $(u_3, u_4)$, which fixes the point $m=(x, 0,0,0)$, so that $g(\partial _{u_2}, \partial _{u_4})(m) = g(\partial _{u_2}, -\partial _{u_4})(m) = 0$. So $g_{24}(x,0, y, 0) = x y \hat{g}_{24}(x,0, y, 0)$ for some smooth function $\hat{g}_{24}$, and
\begin{align*}
	G_{12} &= x ^2 y ^2 \hat{g}_{24} = x ^2 y ^2 \cdot O(1).
\end{align*} 

Around $p$, we have the expansion 
$$\rho = x y (1+ O(d ^2) ).$$
Torus invariance, harmonicity, and $\tilde\tau(p)=0$ give the Taylor expansion
\begin{align*}
	\tilde{\tau } (u) = B x ^2 - B y ^2 + O(d ^4).
\end{align*} 
The gradient identity gives $4B^2=1$. Choose the orientation of $X^\circ$ so that $B=1/2$; reversing the orientation reverses $\tau$. By the definition, we have
\begin{align*}
	\xi  = \frac{1}{2} d ^2 + O(d ^{4}).
\end{align*} 

\end{proof}

\begin{lem}\label{lem-axis}
	Let $E_{+}$ and $E_{-}$ be the boundary rays of $X$, characterized near the corner by $y_p = 0$ and $x_p=0$ respectively as above. Orient both rays away from the corner. Then $\tau$ is strictly increasing on $E_+$ and strictly decreasing on $E_-$. In particular, $\xi $ vanishes only at the corner. At every ordinary axis point, $(\rho , \tau )$ is a local diffeomorphism of closed half-plane neighborhoods.
\end{lem}

\begin{proof}
Set $W= X_1 \wedge X_2$. Then $|W|^2_{g} = \rho ^2 \circ \pi $. On $M ^{\circ}$, we have
\begin{align*}
	\Delta _{g} |W|_g^2 = 4 | \nabla \rho |^2_{h} \circ \pi = 4 | \nabla \tilde{\rho }|^2_g = 4 |\nabla \tilde{\tau }|^2_g.
\end{align*} 
The first and last expressions are smooth on $M$, so their equality holds everywhere. 

Let $F$ be the fixed surface of $X_1$-generated action and $q \in F$ be an ordinary axis point. Then $X_1(q) =0$ and $X_2(q) \neq 0$ is tangent to $F$. Since $W(q) = 0$,
\begin{align*}
	(\Delta _g |W| ^2_g)(q) = 2 |\nabla W|^2_g(q) = 4 |X_2|^2_g(q).
\end{align*}
Indeed, $A:=(\nabla X_1)(q)$ vanishes on $T_qF$ and generates a rotation on $N_qF$ by the Killing equation. Its weight is one in absolute value: a higher weight would give nontrivial finite isotropy at nearby principal points, whereas $X_1$ is primitive of period $2\pi$ and the principal action is free. Choose an orthonormal frame so that $e_1, e_2 \in N_q F$ with $A e_1 = e_2, A e_2 = - e_1$, and $e_3, e_4 \in T_q F$, then $(\nabla _{v} W) (q) = A v \wedge X_2(q)$. Consequently, $|\nabla W|^2_g(q) = |e_2 \wedge X_2(q)|^2 + |e_1 \wedge X_2(q)|^2 = 2|X_2|_g^2(q).$

In particular, $|\nabla \tilde{\tau }|_g(q)>0$ at any ordinary axis point $q$ in $F$. Note that $\nabla \tilde{\tau }(q) \in T_q F$. To see this, let $\sigma $ be the angle $\pi $ rotation generated by $X_1$, then $\sigma (q) = q$ and $d_q \sigma = - \mathrm{Id}$ on $N_q F$. For any $v \in N_q F$, we have $d_q \tilde{\tau }(v) = d_q (\tilde{\tau }\circ \sigma )(v) = d_q \tilde{\tau } (- v)$, which implies that $d_q \tilde{\tau }(v) =0$. Similarly, $\nabla \tilde{\tau }(q) \perp X_2(q)$. This shows that $\tilde{\tau }|_{F}$ descends to $\tau |_{E_{\pm}}$ with $d (\tau |_{E_{\pm}}) \neq 0$ at every ordinary boundary point. So $\tau |_{E_{\pm}}$ is either strictly monotone increasing or decreasing, with the sign determined by the corner expansion. The same calculation with the indices exchanged applies to the other ray.

To show the local diffeomorphism at any ordinary axis point, we use the global equivariant identification $M \simeq \mathbb{C}^2$. The Cartesian coordinates are $z_j = a_j + i b_j, x= |z_1|, y=|z_2|$ and $(x,y) \in [0, \infty) ^2$ are global coordinates on $X$. Note the difference from the normal coordinates $(x_p, y_p)$ we used before. 

Fix $(0, y_0) \in E_{-}$ with $y_0>0$. Choose small enough $\epsilon >0$ and consider the tubular neighborhood $0 \leq x < \epsilon , |y-y_0|< \epsilon $, away from the corner. We can evaluate the metric on the slice $(a_1, b_1, a_2, b_2) = (a_1, 0, y, 0)$ with $|a_1|< \epsilon , |y-y_0|< \epsilon $, where $X_1 = a_1 \partial _{b_1}, X_2 = y \partial _{b_2}$. So
\begin{align*}
	G_{11} = a_1 ^2 g_{b_1 b_1},\ \ G_{12} = a_1 y g_{b_1 b_2},\ \ G_{22} = y ^2 g_{b_2b_2}.
\end{align*} 
Set $\tilde B=(g_{b_1b_1}g_{b_2b_2}-g_{b_1b_2}^2)^{1/2}>0$. On the signed slice the invariant function is $\tilde\rho=|a_1|y\tilde B$. The first-circle half-turn makes $g_{b_1b_1}$ and $g_{b_2b_2}$ even in $a_1$, and $g_{b_1b_2}$ odd in $a_1$. Thus $\tilde B$ is smooth and even, and on the quotient we obtain
\begin{align*}
	\rho = x y B (x ^2, y) \text{ with } B(0, y) >0,\ \ \tau = \tau (x ^2, y),
\end{align*} 
The previous computation shows that $\tau_y(0,y_0)\ne0$. Extend the formulas $xyB(x^2,y)$ and $\tau(x^2,y)$ smoothly to signed $x$; the signed area formula agrees with $\rho$ for $x\ge0$. Their Jacobian at $(0,y_0)$ is $y_0B(0,y_0)\tau_y(0,y_0)\ne0$. The inverse function theorem, restricted to $x\ge0$, proves the asserted local diffeomorphism.

\end{proof}

\begin{lem}\label{lem-green}
The minimal positive Green function of $(M^4, g)$ at $p$, normalized by $- \Delta _{g} G_p = 4 \pi ^2 \delta _p$, satisfies
	\begin{align*}
		G_{p} = \frac{1}{2 \tilde{\xi } }, \ \ \tilde{\xi }:= \xi \circ \pi.
	\end{align*} 
	It is smooth on $M \setminus \{ p\} $ and has the local expansion $G_p \cdot d_g(p, \cdot ) ^2 \to 1$ around $p$.
\end{lem} 
\begin{proof}
Set $U=(2\tilde\xi)^{-1}$. For a smooth function $v(\rho,\tau)$, the conjugacy identities give on $M^\circ$
\[
 \Delta_g(v\circ\pi)
 =\bigl[|d\rho|_h^2(v_{\rho\rho}+\rho^{-1}v_\rho+v_{\tau\tau})\bigr]\circ\pi.
\]
The expression in parentheses vanishes for $v=(2\sqrt{\rho^2+\tau^2})^{-1}$. Moreover, $\tilde\xi^2=|X_1\wedge X_2|_g^2+\tilde\tau^2$ is smooth and positive off $p$, so $U$ extends smoothly and harmonically across every ordinary axis. At $p$, the normal-coordinate expansion gives $U=d^{-2}+O(1)$ with the corresponding differentiated expansion. Its inward flux is therefore $4\pi^2$, proving $-\Delta_gU=4\pi^2\delta_p$.

Dirichlet Green functions on an exhaustion increase to a minimal positive Green function $G_p\le U$. Their difference extends harmonically across $p$ (the pole cancels distributionally) and is nonnegative. Yau's harmonic Liouville theorem \cite{notes-Yau} makes $U-G_p=c$ for some $c\ge0$. If $c>0$, then $U\ge c$, hence $|\tilde\tau|\le\tilde\xi\le(2c)^{-1}$. The same Liouville theorem would make the globally harmonic $\tilde\tau$ constant, contradicting its nonzero quadratic term at $p$. Thus $c=0$ and $G_p=U$.
\end{proof}

\begin{prop}\label{prop-weyl}
		The map
	\begin{align*}
		(\rho , \tau ): X \to \{ (s,t) \in \mathbb{R}^2: s \geq 0\} 
	\end{align*}
	is a proper homeomorphism and a conformal diffeomorphism on the interiors.
\end{prop}
\begin{proof}
	We first show this map is proper. By Li--Yau's Green-function estimate \cite[Theorem~5.2]{notes-LiYau}, we know as $d_{g}(p, \cdot ) \to \infty$, $G_{p} (\cdot ) \to 0$, so that $\xi \circ \pi (\cdot ) \to \infty$. Completeness makes the sublevel sets of $\xi\circ \pi  $ compact on $M$. So $(\rho , \tau ) $ is proper on the quotient.

	Properness and the strict monotonicity of $\tau $ make $\tau $ map the two boundary rays bijectively onto the positive and negative real rays. Thus the boundary map is a homeomorphism. The map $w = \tau + i \rho $ is holomorphic on $X ^{\circ}$. Its nonempty image is open and closed, thus is the whole upper half-plane. 

	Fix an ordinary axis point $q$ and a relatively open neighborhood $U$ on which $w$ is a diffeomorphism of half-plane neighborhoods. There is a neighborhood $V$ of $w(q)$ such that $w^{-1}(V)\subset U$. Otherwise, choose $q_j\notin U$ with $w(q_j)\to w(q)$. Properness gives a convergent subsequence with limit $q_\infty\notin U$; since $\rho(q_\infty)=0$, this contradicts the uniqueness of the boundary preimage of $w(q)$.

Choose an interior value in $V\cap w(U)$. Its only preimage lies in $U$, where the differential is invertible, so its local degree is one. A proper nonconstant holomorphic map has constant total degree, with each local multiplicity positive. Hence its degree is one everywhere, it has no branch points, and it is a conformal diffeomorphism of the interiors. Properness and the already bijective boundary map give the asserted homeomorphism of the closed quotients.
\end{proof}

\subsection{The sharp quotient bound}\label{subsec-quotient-bound}
By Proposition~\ref{prop-weyl}, $(\rho,\tau)$ are coordinates on $X^\circ$.
Put $\delta_2=d\rho^2+d\tau^2$ and $\Delta_2=\Delta_{\delta_2}$.
From this point onward, quotient scalar functions are identified with their
invariant pullbacks when used on $M$. We retain tildes when explicitly
distinguishing a function from its lift; metric subscripts specify all
norms and differential operators.
Applying Colding's sharp estimate \cite[Theorem~3.1]{notes-Colding} to the function $b = G_{p} ^{-\frac{1}{2}} = \sqrt{2 \tilde{\xi }} $ on $(M, g)$, we obtain
\begin{align*}
	|\nabla b|_{g} \leq 1. 
\end{align*} 
On $M ^{\circ}$, we have
\begin{align*}
	|\nabla b|^2_g = \frac{|d \xi |^2_h}{2 \xi } \circ \pi  = \frac{|d \rho |^2_h}{2 \xi }\circ \pi .
\end{align*} 

Note that
\begin{align*}
	d \rho ^2+ d \tau ^2 = |d \rho |^2_{h} h.
\end{align*} 
By Proposition~\ref{prop-weyl}, $|d \rho |_h \neq 0$, so we can write
\begin{align*}
	h = e ^{ 2 \nu } (d \rho ^2+ d \tau ^2) \text{ on } X ^{\circ},
\end{align*} 
where $\nu = - \log |d \rho |_{h}$. Colding's estimate gives
\begin{align*}
	2 \xi \geq |d \rho |^2_h = e ^{-2 \nu }.
\end{align*} 
Integrating the gradient estimate from $p$, where $\tilde\xi=0$, gives $\tilde{\xi } \leq \frac{1}{2} d_{g}(p,\cdot ) ^2$. Set $F:=2 \xi e ^{2 \nu } $.  Then $F \geq 1$ and
\begin{align*}
	h = \frac{F}{2 \xi } (d \rho ^2+ d \tau ^2).
\end{align*} 

\subsection{The harmonic-map equations}\label{subsec-harmonic-map}
The normalized orbit matrix $H=\rho^{-1}G$ takes values in the hyperbolic
plane. Its weighted harmonic-map equation becomes the ordinary harmonic-map
equation after introducing an auxiliary Euclidean three-space.

We first identify the space $\mathcal{P}$ of positive symmetric determinant-one $(2 \times 2)$-matrices with the hyperbolic plane $\mathbb{H}^2$, equipped with the curvature $-1$ metric. For $H \in \mathcal{P}$, we have
\begin{align*}
	T_{H} \mathcal{P} = \{V = V ^{T}: \mathrm{tr}(H ^{-1} V) =0\} ,
\end{align*} 
and $\forall B, C \in T_{H} \mathcal{P}$,
\begin{align*}
	\langle B, C \rangle_{H} := \frac{1}{2} \mathrm{tr}\left( H ^{-1} B H ^{-1} C \right) .
\end{align*}  
Let $D$ denote the Levi-Civita connection of the target metric and also its pullback along $H$. If $i$ denotes a coordinate direction on the domain, $H_i=\partial_iH$, and $V$ is a tangent matrix field along $H$, then
\begin{align*}
	D_{i} V = \partial _i V - \frac{1}{2} (H_i H ^{-1} V + V H ^{-1} H_i).
\end{align*} 
To see the identification, one can take the coordinates
\begin{align*}
	H = \frac{1}{u} \begin{pmatrix} u ^2 + v ^2 & v \\ v & 1 \end{pmatrix} ,\ \ u = H_{22}^{-1} >0, \ \ v= H_{22} ^{-1} H_{12},
\end{align*} 
so that
\begin{align*}
	g_{\mathcal{P}} = \frac{d u^2 + d v ^2}{ u ^2} = g_{\mathbb{H}^2}.
\end{align*} 

In the Weyl coordinates, recall the vertical Ricci equation
\begin{align*}
	0&=\mathrm{div}_{h}\left( \rho H ^{-1} d H \right)\\
	 &=e ^{- 2 \nu }\left( \partial _{\rho }\left( \rho H^{-1} H_{\rho } \right) + \partial _{\tau }\left( \rho H ^{-1} H_{\tau } \right)   \right) .
\end{align*} 
That is, when $\rho >0$,
\begin{align*}
	 H_{\rho \rho } + \rho ^{-1} H_{\rho }- H_{\rho } H ^{-1} H_{\rho } +  H_{\tau \tau } - H_{\tau }H ^{-1} H_{\tau }=0.
\end{align*} 
In terms of the pullback connection, this is
\begin{align*}
	D_{\rho } H_{\rho } + \rho ^{-1} H_{\rho } + D_{\tau } H_{\tau } =0.
\end{align*} 

We introduce an auxiliary $\mathbb{R}^3$ to simplify the arguments. Rotate the half plane $\{(\rho , \tau ): \rho \geq 0\} $ around $\rho =0$, and let $\zeta = (\zeta _1, \zeta _2, \tau )$ denote the Cartesian coordinates. Define $\rho = \sqrt{\zeta _1 ^2+ \zeta _2 ^2 } $, $\mathcal{A}= \{(0,0,\tau ): \tau \in \mathbb{R}\} $ and $\Omega = \mathbb{R}^3 \setminus \mathcal{A}$. On $\Omega $, write $(\zeta _1, \zeta _2) = (\rho \cos\varphi _{\mathrm{aux}}, \rho \sin \varphi _{\mathrm{aux}})$, and lift the map by
\begin{align*}
	H(\zeta ) = H(\rho , \tau ) : \Omega \to \mathbb{H}^2.
\end{align*} 
The auxiliary Euclidean metric and volume are
\begin{align*}
	\delta _{3} = d \rho ^2 + \rho ^2 d \varphi ^2_{\mathrm{aux}} + d \tau ^2, \ \ dV_{\delta _3} = \rho d \rho d \varphi _{\mathrm{aux}} d \tau ,
\end{align*} 
and for a function $f = f(\rho , \tau )$, 
\begin{align*}
	\Delta _{\delta _3} f =  f_{\rho \rho } + \rho ^{-1} f_{\rho } + f_{\tau \tau } \text{ on } \Omega .
\end{align*} 
Write $\Delta_3=\Delta_{\delta_3}$. With the orthonormal frame
$e_1=\partial_\rho$, $e_2=\rho^{-1}\partial_{\varphi_{\mathrm{aux}}}$,
$e_3=\partial_\tau$ on $\Omega$, the tension field of
$H:(\Omega,\delta_3)\to(\mathbb H^2,g_{\mathbb H^2})$ is
\begin{align*}
	\mathcal{T}_{\delta _3}(H) &= \sum_{\alpha }^{} \left( D_{e_{\alpha }} (dH(e_{\alpha }) ) - d H (\nabla ^{\delta _3}_{e_{\alpha }} e_{\alpha }) \right)  \\
				   &= D_{\rho } H_{\rho } + \rho ^{-1} H_{\rho } + D_{\tau } H_{\tau } =0.
\end{align*} 

For base indices $a,b\in\{\rho,\tau\}$, the identities $G=\rho H$ and $\mathrm{tr}(H^{-1}H_a)=0$ give
\begin{align*}
	\mathrm{tr}\left( G ^{-1} G_{a} G ^{-1} G_{b} \right) = 2 (\partial _{a} \log \rho ) (\partial _{b} \log \rho ) +2\langle H_a,H_b\rangle_{\mathbb H^2}. 
\end{align*} 
Thus the horizontal Ricci-flat equation is
\begin{align*}
	\mathrm{Ric}_{h} = \rho ^{-1} \nabla ^2_h \rho - \frac{d \rho \otimes d \rho }{2 \rho ^2} + \frac{1}{2} H ^{*} g_{\mathbb{H}^2}.
\end{align*} 
On the other hand, for $h = e ^{2 \nu } \delta _{2}$, its Christoffel symbols are $\Gamma ^{c}_{a b} = \delta _{a}^{c} \nu _b + \delta ^{c}_{b} \nu _{a} - \delta _{a b} \delta ^{c d} \nu _{d}$, so we have
\begin{align*}
	\mathrm{Ric}_h = - (\Delta _{2} \nu ) \delta _2,\ \ (\nabla ^2_h \rho )_{ab} = \begin{pmatrix} - \nu _{\rho } & - \nu _{\tau } \\ - \nu _{\tau } & \nu _{\rho } \end{pmatrix} _{a b}.
\end{align*} 
Thus we obtain
\begin{align*}
	\nu _{\rho } = - \frac{1}{4 \rho } + \frac{\rho }{4} \left( |H_{\rho }|^2_{\mathbb{H}^2} - |H_{\tau }|^2_{\mathbb{H}^2} \right) ,\ \ \nu _{\tau } = \frac{\rho }{2} \langle H_{\rho }, H_{\tau } \rangle_{\mathbb{H}^2},
\end{align*} 
and
\begin{align*}
	-\Delta _2 \nu  = \rho ^{-1} \nu _{\rho } + \frac{1}{2} |H_{\tau }|^2_{\mathbb{H}^2}.
\end{align*} 

On the auxiliary space, $\xi=|\zeta|$ and $F=2\xi e^{2\nu}\ge1$. The scalar equations are
\begin{align*}
	\Delta _{3} \nu = -\frac{1}{2} |H_{\tau }|^2_{\mathbb{H}^2},\ \ \Delta _{3} \log F= \xi ^{-2} - |H_{\tau }|^2_{\mathbb{H}^2}.
\end{align*} 

\subsection{The sharp Jacobi estimate}\label{subsec-jacobi}
\begin{prop}\label{prop-jacobi}
The norm of $J=H_\tau$, initially defined on $\Omega$, extends to a locally
Lipschitz, weakly subharmonic function on $\mathbb R^3\setminus\{0\}$.
This extension satisfies
\[
 |J|_{\mathbb H^2}\le\xi^{-1}
 \quad\text{on }\mathbb R^3\setminus\{0\}.
\]
\end{prop}
\begin{proof}
Translations in $\tau$ preserve $(\Omega,\delta_3)$, so $H(\zeta_1,\zeta_2,\tau+s)$ is harmonic for every $s$. Its variation field is $J=H_\tau$, a section of $H^*T\mathbb H^2$. Let $E_1=\partial_{\zeta_1}$, $E_2=\partial_{\zeta_2}$, $E_3=\partial_\tau$, and $Y_i=dH(E_i)$. Thus $J=Y_3=dH(E_3)$. Write $D_i=D_{E_i}$. In particular, $D_\tau=D_3$ differentiates a vector field along the target curve $t\mapsto H(\zeta_1,\zeta_2,t)$.

Since $[\partial _{\tau }, E_i] =0$, we have
\begin{align*}
	D_i J = D_{\tau } Y_i,\ \ (D_{\tau }D_{i} - D_{i} D_{\tau }) V= \mathrm{R}_{\mathbb{H}^2}(J, Y_i) V.
\end{align*} 
Since the Cartesian frame is parallel, the harmonic map equation is $\sum_i D_iY_i=0$. Applying $D_\tau$ and using the preceding commutator yields the harmonic-map Jacobi equation
\begin{align*}
	0 &= \sum_{i=1}^{3} (D_iD_i J + R_{\mathbb{H}^2}(J, Y_i) Y_i)\\
	&= \Delta ^{D} J - \sum_{i=1}^{3} \left( |Y_i|^2_{\mathbb{H}^2} J - \langle J, Y_i \rangle_{\mathbb{H}^2} Y_i \right)  . 
\end{align*} 
Thus
\begin{align*}
	\frac{1}{2} \Delta _{3}|J|^2_{\mathbb{H}^2} &= |D J|^2_{\delta _3, \mathbb{H}^2} + \langle \Delta ^{D} J, J \rangle_{\mathbb{H}^2}\\
						    &= |D J|^2_{\delta _3, \mathbb{H}^2} + \sum_{i=1}^{3} \left( |J|^2_{\mathbb{H}^2} |Y_i|^2_{\mathbb{H}^2} - \langle J, Y_i \rangle_{\mathbb{H}^2}^2\right) \\
						    &\geq |D J|^2_{\delta _3, \mathbb{H}^2}\\
						    & \geq |\nabla ^{\delta _3} |J|_{\mathbb{H}^2}|^2_{\delta _3}.
\end{align*} 
Regularization gives
\begin{align*}
	\Delta _{3} \sqrt{|J|^2_{\mathbb{H}^2} + \epsilon ^2} \geq \frac{\epsilon ^2}{(|J|^2_{\mathbb{H}^2} + \epsilon ^2) ^{3 /2} } |D J|^2_{\delta _3, \mathbb{H}^2} \geq 0,
\end{align*} 
taking $\epsilon \to 0$, we have 
$$\Delta _3 |J|_{\mathbb{H}^2} \geq 0 \text{ weakly on } \Omega .$$

We extend this weak subharmonic inequality through the ordinary axis. Fix an axis point with $\tau_0\ne0$, and exchange generators locally if necessary so that $X_1$ collapses. On the signed action slice from Lemma~\ref{lem-axis}, write
\[
 g_{b_2b_2}=\ell(a_1^2,y),\qquad g_{b_1b_2}=a_1k(a_1^2,y).
\]
For $P=x^2$, this gives $G_{22}=y^2\ell(P,y)$ and $G_{12}=Py k(P,y)$. The map
\[
 (P,y)\longmapsto(\rho^2,\tau)
       =(Py^2B(P,y)^2,\tau(P,y))
\]
has Jacobian $y_0^2B(0,y_0)^2\tau_y(0,y_0)\ne0$ at $(0,y_0)$. Its smooth inverse shows that
\[
 G_{22}=A(\rho^2,\tau)>0,\qquad G_{12}=\rho^2C(\rho^2,\tau).
\]
Indeed $G_{12}/\rho^2=k(P,y)/(yB(P,y)^2)$ is smooth there. Thus the target coordinates satisfy
\[
 u=\rho/A,\qquad v=\rho^2C/A,\qquad
 |J|_{\mathbb H^2}^2
   =(A_\tau/A)^2+\rho^2A^2\bigl((C/A)_\tau\bigr)^2.
\]
Consequently $|J|$ is the norm of the smooth three-vector
\[
 \bigl(A_\tau/A,\ \zeta_1A(C/A)_\tau,\ \zeta_2A(C/A)_\tau\bigr),
\]
and is locally Lipschitz through the axis. Test its weak inequality against $\varphi\eta(\rho/\epsilon)$ on a compact set avoiding zero. The extra gradient term is $O(\epsilon)$: $|\nabla^{\delta_3}\eta(\rho/\epsilon)|\le C/\epsilon$, the tube has volume $O(\epsilon^2)$, and $|\nabla^{\delta_3}|J||$ is locally bounded. Passing to the limit proves
\[
 \Delta_3|J|_{\mathbb H^2}\ge0
 \quad\hbox{weakly on }\mathbb R^3\setminus\{0\}.
\]
The square $|J|^2$ is smooth there; its subharmonicity extends by continuity. 

Next, we show that $|J|_{\mathbb{H}^2} \to 0$ uniformly at infinity. Recall that $\Delta _{3} \nu = -\frac{1}{2} |J|^2_{\mathbb{H}^2}$ and $\nu  \geq -\frac{1}{2} \log (2 |\zeta |)$. Note that $e ^{- 2 \nu }\circ \pi  = |\nabla \tilde{\tau }|^2_g$ is smooth, invariant, and strictly positive at ordinary axis points, thus smooth in $(\rho ^2, \tau )$. For the auxiliary radius $r=|\zeta|$, write
\begin{align*}
	\bar{\nu }(r) := \fint_{S^2} \nu (r \omega ) d \omega ,\ \ Q(r):= \int_{1} ^{r} s ^2\overline{ |J|^2_{\mathbb{H}^2}} (s) ds ,\ \ r \geq 1.  
\end{align*} 
Then 
\begin{align*}
	(r ^2 \bar{\nu }'(r) )' = -\frac{1}{2} r ^2 \overline{ |J|^2_{\mathbb{H}^2}} (r).
\end{align*} 
Therefore,
\begin{align*}
	\bar{\nu }(1) - \bar{\nu }(2 r) &= - \bar{\nu }'(1) \left( 1- \frac{1}{2r} \right) + \frac{1}{2 }\int_{1}^{2r} \frac{Q(s)}{s ^2} ds \\
	&\geq - |\bar{\nu }'(1)| + \frac{Q(r)}{4r}.
\end{align*} 
Using $\bar{\nu }(2r) \geq -\frac{1}{2} \log (4r)$, we have
\begin{align*}
	Q(r) \leq C r(1+ \log r).
\end{align*} 
Note that $|J|^2_{\mathbb{H}^2}$ is also subharmonic. For $|\zeta | = r \geq 2$, applying the mean-value inequality on the ball  $B_{\delta _3}(\zeta , r /2)$ gives
\begin{align*}
|J|^2_{\mathbb{H}^2}(\zeta)
&\leq Cr^{-3}
 \int_{B_{\delta_3}(\zeta,r/2)}|J|^2_{\mathbb{H}^2}\,dV_{\delta_3}\\
&\leq Cr^{-3}Q(2r)
\leq C r ^{-2}(1+\log r).
\end{align*}

Next, we show that around the corner, $|J|_{\mathbb{H}^2} = \xi ^{-1} + O(1) $ uniformly in direction. In the local normal coordinates $(x_p,y_p)$ of Lemma~\ref{lem-conjugate}, for $s_1= x_p ^2, s_2= y_p^2, d ^2 = s_1+ s_2$, we have
\begin{align*}
	G_{11} &= s_1 A_1 ,\ \ G_{22}= s_2 A_2, \ \ G_{12} = s_1 s_2 C,
\end{align*} 
where $A_1,A_2,C$ are smooth functions of $(s_1,s_2)$, $A_i=1+O(d^2)$, and $C=O(1)$. Their first derivatives with respect to $(s_1,s_2)$ are bounded. Smooth torus invariance also gives
\[
 \rho^2=s_1s_2(A_1A_2-s_1s_2C^2),\qquad
 A_1A_2-s_1s_2C^2=1+O(d^2),
\]
\[
 \begin{aligned}
 \tau&=(s_1-s_2)/2+E(s_1,s_2),&\quad \xi&=d^2/2+O(d^4),\\
 E&=O(d^4),& d_{(s_1,s_2)}E&=O(d^2).
 \end{aligned}
\]
The two rows of the Jacobian of $(\rho^2,\tau)$ are
\[
 \bigl(s_2(1+O(d^2)),s_1(1+O(d^2))\bigr),\qquad
 \bigl(1/2+O(d^2),-1/2+O(d^2)\bigr).
\]
Its determinant is $-d^2/2+O(d^4)$. Inverting at fixed $\rho$ therefore gives, uniformly up to either axis,
\[
 \partial_\tau s_1=2s_1/d^2+O(s_1),\qquad
 \partial_\tau s_2=-2s_2/d^2+O(s_2).
\]
Recall the hyperbolic coordinates $(u,v)$ defined above. Then
\begin{align*}
	u = \frac{\rho }{s_2 A_2} ,\ \ v = \frac{s_1 C}{A_2}.
\end{align*} 
So
\begin{align*}
	\frac{u_{\tau }}{u} = - \frac{\partial _{\tau } s_2}{s_2} - \partial _{\tau } \log A_2= \frac{2}{d ^2} + O(1) = \frac{1}{\xi }+ O(1).
\end{align*} 
Also $\partial_\tau v=O(s_1/d^2)$ and $\rho\asymp\sqrt{s_1s_2}$, so
\[
 \left|\frac{v_\tau}{u}\right|
 \le C\frac{\sqrt{s_1s_2}}{s_1+s_2}\le C.
\]

Thus
\begin{align*}
	|J|_{\mathbb{H}^2} = \frac{1}{\xi } + O(1) \text{ as } d \to 0.
\end{align*} 
For $\epsilon>0$, compare the subharmonic $|J|$ with $(1+\epsilon)/\xi+\epsilon$ on an auxiliary Euclidean annulus. The inner boundary inequality holds for sufficiently small inner radius by the uniform corner estimate; the outer inequality holds for sufficiently large outer radius by decay at infinity. The ordinary axes are already included in the weak inequality and create no boundary term. Exhausting and then letting $\epsilon\downarrow0$ gives
\begin{align*}
	|J|_{\mathbb{H}^2} \leq \frac{1}{\xi } \text{ on } \mathbb{R}^3 \setminus \{0\} .
\end{align*} 
So
\begin{align*}
	\Delta _{3} \log F = \xi ^{-2} - |J|^2_{\mathbb{H}^2} \geq 0 \text{ on } \mathbb{R}^3 \setminus \{0\} .
\end{align*} 

\end{proof}

\subsection{The conformal factor and a volume lower bound}\label{subsec-volume}
\begin{prop}\label{prop-conformal-volume}
The conformal factor satisfies
\[
 1\le F\le C(1+\xi),\qquad
 |\nabla^{\delta_3}\log F|_{\delta_3}\le C/\xi\quad(\xi>0).
\]
Moreover, $\Vol_g(B_g(p,R))\ge cR^3$ for every $R\ge1$.
\end{prop}
\begin{proof}
Put $L=\log F$ and $f=\xi^{-2}-|J|_{\mathbb H^2}^2$. Near zero,
\[
 e^{-2\nu}=|\nabla^g\tilde\tau|_g^2=d^2+O(d^4)=2\xi+O(\xi^2),
 \qquad L=O(\xi),\qquad f=O(\xi^{-1}).
\]
Set $L(0)=0$. The punctured equation $\Delta_3L=f$ holds distributionally on all of $\mathbb R^3$: test against $\chi_\epsilon\varphi$, where $\chi_\epsilon$ vanishes on the ball of radius $\epsilon$ and equals one outside radius $2\epsilon$. The bounds $|\nabla^{\delta_3}\chi_\epsilon|\le C/\epsilon$, $|\Delta_3\chi_\epsilon|\le C/\epsilon^2$, and $L=O(\xi)$ make the cutoff errors $O(\epsilon^2)$. The omitted source integral tends to zero as well. Hence
\[
 \Delta_3L=f\quad\hbox{on }\mathbb R^3,\qquad
 L\ge0,\qquad 0\le f\le\xi^{-2}.
\]
Define
\[
 m(r)=\frac1{4\pi}\int_{S^2}L(r\omega)\,d\omega,
 \qquad \mathcal F(r)=e^{m(r)}.
\]
Spherical averaging, with no point source at zero, gives
\[
 0\le m'(r)=\frac1{4\pi r^2}\int_{B_{\delta_3}(0,r)}f\,dV_{\delta_3}
       \le\frac1r.
\]
Thus $m(r)\le\log r+C$ for $r\ge1$. The mean-value inequality for nonnegative subharmonic $L$ now gives the pointwise bound $L(\zeta)\le C\log(2+|\zeta|)$.

The normalized Newton representation is
\begin{equation}\label{notes-newton}
 L(\zeta)=\frac1{4\pi}\int_{\mathbb R^3}
 \left(\frac1{|\eta|}-\frac1{|\zeta-\eta|}\right)f(\eta)\,d\eta.
\end{equation}
Denote the right side by $P(\zeta)$. Splitting at $|\eta|=2|\zeta|$ and at $|\eta|=1$, using $f=O(|\eta|^{-1})$ near zero and $f\le|\eta|^{-2}$ elsewhere, proves absolute convergence,
\[
 P(\zeta)=O\bigl(|\zeta|(1+|\log|\zeta||)\bigr)\quad(\zeta\to0),
 \qquad |P(\zeta)|\le C\log(2+|\zeta|)\quad(|\zeta|\ge1).
\]
The kernel identity gives $\Delta_3P=f$. Hence $L-P$ is an entire harmonic function of logarithmic growth, so it is constant by the interior harmonic gradient estimate. It vanishes at zero, proving \eqref{notes-newton}.

Differentiating for $\zeta\ne0$ yields
\[
 |\nabla^{\delta_3}L|_{\delta_3}(\zeta)
 \le\frac1{4\pi}\int_{\mathbb R^3}
       \frac{d\eta}{|\eta|^2|\zeta-\eta|^2}
 \le\frac C{|\zeta|}.
\]
The last integral is finite near both singularities and at infinity; scaling gives $|\zeta|^{-1}$. Integrating this gradient bound along spherical arcs and radial segments gives
\[
 \operatorname{osc}_{\{r/2\le|\zeta|\le r\}}L\le C,
 \qquad C^{-1}\mathcal F(r)\le F(r\omega)\le C\mathcal F(r).
\]
Moreover $0\le m'\le1/r$ implies
\[
 \mathcal F\text{ is nondecreasing},\qquad
 \mathcal F(2r)\le2\mathcal F(r),\qquad
 1\le\mathcal F(r)\le C(1+r).
\]
Consequently
\[
 1\le F\le C(1+\xi).
\]

Finally, we estimate the volume growth of $g$. Write the quotient coordinate $(\rho ,\tau )$ as $\rho = \xi \sin \theta , \tau = \xi \cos \theta $. Then
\begin{align*}
	h= \frac{F}{2 \xi }(d \xi ^2+ \xi ^2 d \theta ^2).
\end{align*} 
The radial segment from the corner to $(r, \theta )$ has length
\begin{align*}
	\ell_{\theta }(r) = \int_{0} ^{r} \sqrt{ \frac{F(\xi , \theta )}{2 \xi } }d \xi .
\end{align*} 
The bound of $F$ implies that
\begin{align*}
	\sqrt{2 r}  \leq \ell_{\theta }(r) \leq C (\sqrt{r} + r).
\end{align*} 
Since quotient distance from $\pi (p)$ equals $g$-distance from $p$, for any $q$ over $(r, \theta )$, we have
\begin{align*}
	d_{g}(p, q) \leq \ell_{\theta }(r).
\end{align*} 
Since $\rho = \sqrt{\mathrm{det}G} $, we have
\begin{align*}
	\pi _{*}(d V_g) = 4 \pi ^2 \rho d V_h = 2 \pi ^2 F(\xi , \theta ) \xi \sin \theta d \xi d \theta .
\end{align*} 
Fix a small constant $\epsilon >0$. For all large radius $R$, for all $\theta $, we have
\begin{align*}
	\ell_{\theta }(\epsilon R) \leq \frac{1}{2}R, \ \ \ell_{\theta }(R'_{\theta }) = R,
\end{align*} 
for some $\epsilon R< R'_{\theta }$. So the domain over the radial region $\xi < R'_{\theta }$ is inside $B_g(p, R)$. Note that
\begin{align*}
	\frac{1}{2}R \leq \ell_{\theta }(R'_{\theta }) - \ell_{\theta }(\epsilon R) &= \int_{\epsilon R}^{R'_{\theta }} \sqrt{\frac{F(\xi ,\theta )}{2 \xi }} d \xi \leq \left( \frac{1}{2 \epsilon R} \int_{\epsilon R}^{R'_{\theta }} \xi F(\xi ,\theta ) d \xi \right) ^{\frac{1}{2}}.
\end{align*} 
Therefore,
\begin{align*}
	\mathrm{Vol}_g(B_g(p,R) ) &\geq 2 \pi ^2 \int_{0}^{\pi } \sin \theta  \int_{\epsilon R} ^{R'_{\theta }} \xi F(\xi ,\theta ) d \xi d \theta \\
				  &\geq 2 \pi ^2 \epsilon R ^3.
\end{align*} 
Reducing $c$ to cover bounded radii gives
\begin{align*}
	\mathrm{Vol}_g(B_g(p,R) ) \geq c R^3,\ \ \forall R \geq 1.
\end{align*} 

\end{proof}

\subsection{The asymptotic dichotomy}\label{subsec-asymptotic}

\begin{prop}\label{notes-charge}
Either $g$ is Euclidean, or, uniformly as $\xi\to\infty$,
$$
 \xi |J|_{\mathbb{H}^2}\to 0,\qquad
 \xi|\nabla^{\delta_3}\nu|_{\delta_3}\to 0,
 \qquad F= 2 \xi e ^{2 \nu }\to \infty.
$$
In the latter case, after target isometries, every subsequential interior limit
of $H(t\rho,t\tau)$ is a unit-speed hyperbolic geodesic parametrized by $\log\rho$.
\end{prop}
\begin{proof}
Set $L= \log F$, $a = |J|_{\mathbb{H}^2}$ and $q(r)=r\bar a(r)$, where the bar denotes the spherical average in
$\mathbb R^3$. Since $\xi =|\zeta |$, $a(\zeta ) \leq \frac{1}{|\zeta |}$ and $a(\zeta ) = \frac{1}{|\zeta |} + O(1)$ as $\zeta \to 0$, so $0\le q\le1$ and $q(0+)=1$.
For a nonnegative $\varphi\in C_c^\infty((0,\infty))$,
\[
	\Delta_3\left(\frac{\varphi(r)}{r}\right)
	=\frac{\varphi''(r)}{r},
 \qquad
 4\pi\int_0^\infty q(r)\varphi''(r)\,dr = \int_{\mathbb{R}^3} r ^{-1} \varphi (\Delta _3 a) dV_{\delta _3} \ge0.
\]
Thus $q$ is convex. A convex function bounded above on $(0,\infty)$ is
nonincreasing, so $\lim_{r\to \infty}  q(r)= c$ for some $c\in[0,1]$.

We first prove that, for every $1\le p<\infty$, as $t\to \infty$,
\begin{equation}\label{notes-charge-Lp}
 a_t(\zeta):=t a(t\zeta)\to  c/|\zeta|
 \quad\hbox{in }L^p_{\mathrm{loc}}(\mathbb R^3\setminus\{0\}).
\end{equation}
Indeed, $a_t$ is nonnegative, subharmonic, locally Lipschitz on $\mathbb{R}^3 \setminus \{0\} $, and bounded by $|\zeta|^{-1}$.
For a cutoff $\chi $ supported away from zero, integration by parts gives
\[
 \int\chi^2|\nabla^{\delta_3}a_t|_{\delta_3}^2
       \le4\int a_t^2|\nabla^{\delta_3}\chi|_{\delta_3}^2.
\]
Rellich compactness therefore, up to a subsequence $t_j \to \infty$, $a_{t_j} \to a_\infty$ in strong local $L^2$-topology, with $0\le a_\infty \le|\zeta|^{-1}$ and $\Delta_3 a_\infty \ge0$. Since
$r\overline{a_{t_j}}(r)=q(t_j r)\to c$ as $t_j\to \infty$, by Jensen’s inequality, we know $r \overline{a_\infty} (r) = c$ for almost every
$r$. Together with mean value inequality of $\Delta _3 a_\infty \geq 0$, this proves the mean value equality, thus $\Delta _3 a_\infty =0$. The Kelvin
transform
\[
	\widehat {a_\infty}(\zeta)=|\zeta|^{-1}a_\infty(\zeta/|\zeta|^2)
\]
is harmonic and bounded between zero and one on $\mathbb R^3\setminus\{0\}$, and thus is a constant by the Liouville theorem. 
Using the spherical mean identifies again, we obtain $a_\infty = c/|\zeta|$, independent of the choice of subsequence. The local uniform bound upgrades convergence to every
finite $L^p$. It also gives $a_t^2\to c^2/|\zeta|^2$ in every finite local $L^p$.

Fix $\zeta_0=(1,0,0)$, so $(\rho,\tau)=(1,0)$, and set
$L_t(\zeta)=L(t\zeta)-L(t\zeta_0) = \log \frac{|\zeta |}{|\zeta _0|} + 2 \nu (t \zeta ) - 2 \nu (t \zeta _0)$. The gradient estimate and
\[
 \Delta_3L_t=|\zeta|^{-2}-a_t^2
\]
give local $W^{2,p}$ bounds on every annulus, including its rotation axis.
Subsequential limits satisfy
\[
 \Delta_3L_\infty=(1-c^2)|\zeta|^{-2},\qquad
 |\nabla^{\delta_3}L_\infty|_{\delta_3}\le C/|\zeta|,
 \qquad L_\infty(\zeta_0)=0.
\]
The harmonic function $L_\infty-(1-c^2)\log|\zeta|$ has at most logarithmic
growth both at zero and infinity. So it can be extended to a global harmonic function and thus a constant zero by Liouville theorem. Thus, for every $0<\alpha<1$,
\begin{equation}\label{notes-scalar-limit}
 L_t\to (1-c^2)\log|\zeta |,\ \ 
 \nu(t\zeta)-\nu(t\zeta_0)\to -\tfrac12c^2\log|\zeta |
 \ \text{ in }C^{1,\alpha}_{\mathrm{loc}}(\mathbb R^3\setminus\{0\}).
\end{equation}

Next, we show that $0<c<1$ is impossible. For each $t>0$, let $I_t$ be a target isometry that
sends $H(t\zeta_0)$ to $(1,0)\in\mathbb H^2$. More precisely, if $H(t \zeta _0) = (U_t, V_t)$, then we can take $I_t(u,v) = \left( \frac{u}{U_t}, \frac{v- V_t}{U_t} \right) $. Define $\mathcal H_t(\zeta)=I_t(H(t\zeta))$. Then
\begin{align*}
	|\partial_\tau\mathcal H_t|_{\mathbb H^2}^2(\zeta ) &=t^2 |H_{\tau }|^2_{\mathbb{H}^2}(t \zeta ) = a_t^2(\zeta ),\\
	|\partial_\rho\mathcal H_t|_{\mathbb H^2}^2(\zeta ) &= t ^2| H_{\rho }(t \rho , t \tau )|^2_{\mathbb{H}^2}							 =\rho^{-2}+4\rho^{-1}\partial_\rho\nu^t(\zeta)+a_t^2(\zeta ),\\
	\langle \partial _\rho \mathcal{H}_t, \partial _\tau \mathcal{H}_t \rangle_{\mathbb{H}^2}(\zeta ) &= t ^2\langle H_{\rho }, H_{\tau } \rangle_{\mathbb{H}^2}(t \rho , t \tau  )= 2 t \rho ^{-1} \nu _{\tau }(t \rho , t \tau ) = 2 \rho ^{-1} \partial _{\tau } \nu ^{t} (\zeta ),
\end{align*} 
where $\nu ^t (\zeta ) := \nu (t \zeta ) - \nu (t \zeta _0)$.
Interior elliptic estimates
for the harmonic-map equation then yield smooth subsequential convergence $\mathcal{H}_{t_j}(\zeta ) \to \mathcal{H}_\infty(\zeta )$
on $\rho>0$. By \eqref{notes-charge-Lp}--\eqref{notes-scalar-limit},
\[
 |\partial _\tau \mathcal{H}_{\infty}|_{\mathbb H^2}^2=\frac{c^2}{|\zeta|^2},\qquad
 |\partial _\rho \mathcal{H}_{\infty}|_{\mathbb H^2}^2=\frac1{\rho^2}-\frac{c^2}{|\zeta|^2},\qquad
 \langle \partial _\rho \mathcal{H}_{\infty}, \partial _\tau \mathcal{H}_{\infty}\rangle_{\mathbb H^2}
       =-\frac{c^2\tau}{\rho |\zeta|^2}.
\]
Introduce only for this computation $\sigma=\log\rho$ and
$\upsilon=\operatorname{arsinh}(\tau/\rho)$. Substitution gives
\[
 \mathcal{H}_\infty^*g_{\mathbb H^2}=(1-c^2)d\sigma^2+c^2d\upsilon^2 \text{ on } X ^{\circ}.
\]
For $0<c<1$ this is a positive definite flat metric, so $\mathcal{H}_\infty$ is a
local diffeomorphism whose pullback metric has curvature zero, a contradiction with the target curvature $-1$. Thus $c=0$ or $1$.

If $c=1$, monotonicity and $q(0+)=1$ give $q\equiv1$. The pointwise inequality
$a\le\xi^{-1}$ then gives equality everywhere. Hence $\Delta_3L=0$ on all
of $\mathbb R^3$; $L\ge0$ and $L(0)=0$ imply $L=0$. Thus
$F \equiv 1$, and the equality case of Colding's
gradient estimate gives the Euclidean metric
\cite[Theorem~3.1]{notes-Colding}.

In the nonflat case $c=0$, \eqref{notes-charge-Lp} gives
\begin{align*}
	a_t \to 0 \text{ in } L^{2}_{\mathrm{loc}}(\mathbb{R}^3 \setminus \{0\} ).
\end{align*} 
The mean-value inequality for the subharmonic
$a_t$ implies that $a_t(\zeta ) = t a(t \zeta ) \to 0$ locally uniformly on $\mathbb{R}^3\setminus \{0\} $. In particular, $\xi a \to 0$ uniformly as $\xi \to \infty$.  Equation
\eqref{notes-scalar-limit} gives $ \nu ^{t}(\zeta ) \to 0 $ in $C^{1,\alpha }_{\mathrm{loc}}(\mathbb{R}^3 \setminus \{0\} )$, so $|\nabla ^{\delta _3} \nu ^{t}|(\zeta ) = t|\nabla ^{\delta _3} \nu |(t \zeta ) \to 0$ locally uniformly. In particular, $\xi|\nabla^{\delta_3}\nu|_{\delta_3}\to0$ as $\xi \to \infty$.
Integration along rays gives $\nu=o(\log\xi)$ uniformly. Therefore
$\log F= \log 2 \xi + 2 \nu =\log 2\xi+o(\log\xi) \to \infty$ uniformly as $\xi \to\infty$. Finally the interior limit
has $\partial _{\tau }\mathcal{H}_{\infty}=0$ and $|\partial _{\rho } \mathcal{H}_{\infty}|_{\mathbb H^2}=1/\rho$;
its harmonic-map equation becomes the unit-speed geodesic equation in
$\sigma=\log\rho$.
\end{proof}

We give a consequence of this asymptotic behavior. Using $\xi |\nabla ^{\delta _2} \nu |_{\delta _2} \to 0$ uniformly as $\xi \to \infty$, we know $\sup_{\xi=s}|\nu(\rho,\tau)-\nu(s,0)|=o(1)$ as $s\to\infty$. Thus $e ^{\nu (\rho , \tau )} = (1+ o(1) ) e ^{\nu (s,0)}$ uniformly, and
\begin{align*}
	d_g(p, q) = (1+ o(1) ) r_s \text{ uniformly for } \tilde{\xi }(q) = s, \ \ \text{ as } s\to \infty,
\end{align*} 
where $r_s := s e ^{\nu (s, 0)}$.

\subsection{Adapted Killing fields and uniform estimates}
Assume $g$ is nonflat. We choose a constant torus Killing field on each
large annulus and estimate its logarithmic derivative through both axes. Recall $r_s = s e ^{\nu (s,0)}$.
\begin{lem}\label{notes-killing}
For every sufficiently large $s$, there is a constant $\omega^s$ with
$1/2\le|\omega^s|\le2$ such that the globally defined Killing field
$$
 K_s=-\omega^sX_1+X_2
$$
satisfies
$$
	\frac{|\nabla^gK_s|_g}{|K_s|_g}\le\frac{C}{r_s}
	\quad\text{on }\{s\le \tilde{\xi}\le2s\}.
$$
\end{lem}
\begin{proof}
We first control the collapsing circle near each axis, then use a scalar
elliptic equation to choose one constant coefficient vector for both axes.
Set
\[
 \sigma_s=s\sup_{s/16\le\xi\le16s,\ \rho>0}
 \bigl(|J|_{\mathbb H^2}+|\nabla^{\delta_3}\nu|_{\delta_3}\bigr)
 \to 0.
\]
Path integration on this larger annulus gives
\begin{equation}\label{notes-annular-conformal}
 e^{2\nu(\rho,\tau)}=e^{2\nu(s,0)}(1+O(\sigma_s)).
\end{equation}
The negative axis $\tau<0$ is where $X_1$ collapses.

Consider the region $0<\rho<3s$, $-4s<\tau<-s/4$, and put
$k=G_{11}/\rho^2$. Recall that in Lemma \ref{lem-axis}, for $q$ over $(0, \tau )$, we have $|X_1|(q) =0$,
\begin{align*}
	|\nabla \tilde{\tau }|^2_g(q) = |X_2|^2_g(q) \neq 0.
\end{align*} 
On the quotient, we obtain $|d \tau |^2_{h}( 0, \tau ) = e ^{-2 \nu (0, \tau )} = G_{22}(0, \tau )$. On the other hand, $k ^{-1} = G_{22} - \frac{G_{12}^2}{G_{11}}$, so $k ^{-1}(0, \tau ) = G_{22}(0, \tau )$. 
Thus
\[
 k(0,\tau)=e^{2\nu(0,\tau)},
\]
and $k$ is smooth in $(\rho^2,\tau)$. The Weyl equation gives
\[
 \bigl||H_\rho|_{\mathbb H^2}^2-\rho^{-2}\bigr|
 \le\frac{C\sigma_s}{s\rho},\qquad
 |dH|_{\delta_3,\mathbb H^2}\le C/\rho.
\]
Recall that $\Omega = \mathbb{R}^3 \setminus \mathcal{A}$  is the auxiliary Euclidean space. Fix $\zeta_0\in\Omega$ with $s/8\le\xi(\zeta_0)\le8s$, and set $\rho_0=\rho(\zeta_0)$ and $R=\rho_0/8$. Then $B_{\delta_3}(\zeta_0,2R)$ avoids the axis, and $\rho$ is comparable to $\rho_0$ throughout this ball. Similarly as in previous section, we normalize the target by an isometry. More precisely, if $H(\zeta _0) = (u_0, v_0)$, then take $I(u, v) = \left( \frac{u}{u_0}, \frac{v-v_0}{u_0} \right) $ so that $I(H(\zeta _0) ) = (1,0) \in \mathbb{H}^2$. Define, for $y \in B_2(0) \subset \mathbb{R}^3$, 
$$\tilde{H}(y) = I(H(\zeta _0 + R y) ).$$ 
Then $\tilde H(0)=(1,0)$ and the gradient estimate gives
\begin{align*}
	|d \tilde{H}|_{\delta _3, \mathbb{H}^2} \leq C.
\end{align*} 
The harmonic-map equation gives uniform interior bounds for the derivatives of $\tilde{H}$, by elliptic bootstrapping. Note that $\tilde{J}(y) = d I (J(\zeta _0 + R y) )$ satisfies the Jacobi equation, whose coefficients are controlled by the bounds of $\tilde{H}, d \tilde{H}$. So interior estimates give
\begin{align*}
	|\tilde{D} \tilde{J}|_{\delta _3,\mathbb{H}^2}(0) \leq C \sup_{B_1(0)} |\tilde{J}|_{\mathbb{H}^2}. 
\end{align*} 
Therefore,
\begin{equation}\label{notes-J-derivative}
 |D_\tau J|_{\mathbb H^2}(\rho , \tau ) \le \frac{C}{\rho}\sup|J|_{\mathbb H^2}
 \le\frac{C\sigma_s}{s\rho} \ \ \text{ on } \{s /8 \leq \xi \leq 8 s\}.
\end{equation}

Use the upper-half-plane chart $u(H)=H_{11}^{-1}$, $v(H)=H_{12}/H_{11}$, obtained from the earlier chart by interchanging the two matrix indices. Then $H_{11}=\rho k=1/u(H)$. Note that $\operatorname{Hess}_{\mathbb H^2}\frac{1}{u}= \frac{1}{u} g_{\mathbb H^2}$
and $|d \frac{1}{u}|_{\mathbb H^2}=\frac{1}{u}$. For simplicity, set $\mathcal{T} = \frac{1}{u}$ and $T = H_{11} = \mathcal{T}(H)$.
The harmonic-map equation yields
\[
 T_{\rho\rho}+\rho^{-1}T_\rho
 =|H_\rho|_{\mathbb H^2}^2T-(d \mathcal{T})_H(D_\tau J).
\]
Substitution of $T=\rho k$ and \eqref{notes-J-derivative} gives
\[
 \left|\partial_\rho(\rho^3k_\rho)\right|
 \le \frac{C\sigma_s}{s}\rho^2k.
\]
Fix $\tau$. Since $\rho^3k_\rho\to0$ at the axis, taking integrations give
\[
 |k_\rho(\rho,\tau)|\le\frac{C\sigma_s}{s}
       \sup_{0\le r\le\rho}k(r,\tau),\qquad
 |k(\rho,\tau)-k(0,\tau)|\le\frac{C\sigma_s\rho}{s}
       \sup_{0\le r\le\rho}k(r,\tau).
\]
For $\rho\le3s$, absorb the last supremum term when $s$ is large to obtain 
\begin{align*}
	|\partial _{\rho } \log k(\rho,\tau)|\le\frac{C\sigma_s}{s},\qquad |k(\rho , \tau ) - k(0, \tau )| \leq \frac{C \sigma _s \rho }{s} k(0, \tau ).
\end{align*} 
Also 
$$|\partial_\tau\log k|=|\partial_\tau\log T| = \frac{|(H_{11})_{\tau }|}{|H_{11}|}\le|J|_{\mathbb H^2}.$$
Together with \eqref{notes-annular-conformal} and $k(0, \tau ) = e ^{2 \nu (0, \tau )}$, these facts give
\begin{equation}\label{notes-k-control}
 k=e^{2\nu(s,0)}(1+O(\sigma_s)),\ \ 
 |\nabla^{\delta_2}\log k|_{\delta_2}\le C\sigma_s/s,
\end{equation}
where $ 0 \leq \rho \leq 3 s,\ -4s \leq \tau \leq - s /4$, for all uniformly large $s$.


We now write
\begin{equation}\label{notes-axis-gram}
 G_{ij}d\theta^i d\theta^j
 =k\rho^2(d\theta^1+\omega d\theta^2)^2+k^{-1}(d\theta^2)^2,
 \qquad \omega=G_{12}/G_{11}.
\end{equation}
In Lemma \ref{lem-axis}, we have shown that both $k$ and $\omega$ are smooth in $(\rho^2,\tau)$, aross $\rho =0$.
In the target hyperbolic chart $(u,v)$, $U:=u(H)=(k\rho)^{-1}$, $V:=v(H)=\omega$. The $V$-component of the harmonic map equation gives
$\Delta_3V-2U^{-1}\langle dU,dV\rangle_{\delta_3}=0$. After substituting \(U=(k\rho)^{-1}\) and \(V=\omega\), the equation becomes
\[
\partial_\rho^2\omega+\frac3\rho\partial_\rho\omega+\partial_\tau^2\omega
+2\left(\frac{k_\rho}{k}\partial_\rho\omega+
\frac{k_\tau}{k}\partial_\tau\omega\right)=0.
\] 
To simplify the arguments, we introduce the auxiliary Euclidean space  $(\mathbb{R}^5, \delta _5)$ with coordinates $(z,\tau)\in\mathbb R^4\times\mathbb R$ and identify 
$\rho=|z|$ so that $\Delta_5=\partial_\rho^2+3\rho^{-1}\partial_\rho+\partial_\tau^2$. Then we can rewrite above equation as 
\begin{equation}\label{notes-five-equation}
 \operatorname{div}_{\delta_5}(k^2\nabla^{\delta_5}\omega)=0.
\end{equation}
Smoothness in $\rho^2$
makes this equation regular across $z=0$.

The target hyperbolic metric gives
\[
 |H_\rho|_{\mathbb H^2}^2=(\rho^{-1}+k_\rho/k)^2+(k\rho)^2(\partial_\rho\omega)^2,
 \qquad
 |H_\tau|_{\mathbb H^2}^2=(k_\tau/k)^2+(k\rho)^2(\partial_\tau\omega)^2.
\]
Using the Weyl equation gives
\begin{align*}
(k\rho)^2(\partial_\rho\omega)^2
=\frac{4\nu_\rho-2\partial_\rho\log k}{\rho}
+|J|_{\mathbb H^2}^2-(\partial_\rho\log k)^2.
\end{align*} 
Using \eqref{notes-k-control} and only
$\sigma_s\le1$, we obtain
\[
 |\nabla^{\delta_5}\omega|_{\delta_5}^2
 \le Ce^{-4\nu(s,0)}\left(\frac1{s\rho^3}+\frac1{s^2\rho^2}\right).
\]
On the cylinder $\mathcal C_s=\{|z|<3s,\ -4s<\tau<-s/4\}$, integration
with $dV_{\delta_5}=|\mathbb{S}^3|\rho^3d\rho\,d\tau$ therefore gives
\[
 \int_{\mathcal C_s}|\nabla^{\delta_5}\omega|_{\delta_5}^2dV_{\delta_5}
 \le Cs e^{-4\nu(s,0)}=Cs^5/r_s^4.
\]
By \eqref{notes-k-control}, the coefficient $e^{-4\nu(s,0)}k^2$ is bounded
above and below by positive constants, and its $\delta_5$-gradient is bounded
by $C/s$. We apply the interior gradient estimate to
\eqref{notes-five-equation} after multiplying by the constant $e^{-4\nu(s,0)}$.
Consider the smaller band and its lift
\[
 \mathcal B_s^-=\{0\le\rho\le2s,\ -3s\le\tau\le-s/2\},\qquad
 \widehat{\mathcal B}_s^-=\{|z|\le2s,\ -3s\le\tau\le-s/2\}.
\]
Every point $q\in\widehat{\mathcal B}_s^-$ has Euclidean distance at least
$s/4$ from $\partial\mathcal C_s$, so the Euclidean ball
$B_{s / 8}(q)$ lies in $\mathcal C_s$. Then the interior scalar gradient estimate and Poincar\'e's
inequality give
\[
 |\nabla^{\delta_5}\omega|_{\delta_5}(q)
 \le \frac{C}{s}
 \left(\fint_{B_{s /8}(q)}|\omega-\bar\omega|^2dV_{\delta_5}\right)^{1/2}
 \le C\left(\fint_{B_{s /8}(q)}
 |\nabla^{\delta_5}\omega|_{\delta_5}^2dV_{\delta_5}\right)^{1/2}.
\]
Thus
\begin{equation}
 \begin{aligned}
 \sup_{\widehat{\mathcal B}_s^-}|\nabla^{\delta_5}\omega|_{\delta_5}
 &\le C\left(s^{-5}\int_{\mathcal C_s}
 |\nabla^{\delta_5}\omega|_{\delta_5}^2dV_{\delta_5}\right)^{1/2}
 \le Cr_s^{-2}, 
 \end{aligned}
\end{equation}
which implies that
\begin{align}\label{notes-twist-gradient}
	\sup_{\mathcal B_s^-}|d \omega|_{\delta _2}&\le Cr_s^{-2}.
\end{align} 

On the reflected band $\mathcal B_s^+=\{0\le\rho\le2s,\ s/2\le\tau\le3s\}$,
interchange $X_1,X_2$ and use
$\widetilde k=G_{22}/\rho^2$, $\widetilde\omega=G_{12}/G_{22}$.

Set $\omega^s=\omega(s,0)=G_{12}(s,0)/G_{11}(s,0)$.
The bound for $k$ at $(s,-s)$ controls $G_{11}$, and its positive-axis
counterpart at $(s,s)$ controls $G_{22}$. Since
\[
 |\partial_\tau\log G_{ii}|\le |J|_{\mathbb H^2}\le\sigma_s/s
 \quad\text{on }\rho=s,\ -s\le\tau\le s,
\]
integration gives, on this whole segment,
\[
 G_{11}=r_s^2(1+O(\sigma_s)),\qquad
 G_{22}=r_s^2(1+O(\sigma_s)).
\]
At $(s,0)$, $\det G=s^2$, and $s/r_s^2=2/F(s,0)\to0$ as $s\to \infty$. Hence
$(\omega^s)^2=G_{22}/G_{11}-s^2/G_{11}^2\to1$, so
$1/2\le|\omega^s|\le2$ for large $s$.
The estimate of $|J|_{\mathbb{H}^2}$ also gives 
$$|\partial_\tau\omega|\le C/r_s^2\ \  \text{ on } \rho =s,\ |\tau | \leq s.$$
Combining this with \eqref{notes-twist-gradient} yields
\begin{equation}\label{notes-twist-match}
 |\omega-\omega^s|\le Cs/r_s^2
 \quad\text{on }\mathcal B_s^-.
\end{equation}
For $\mathcal B_s^+$, use the exact identity
\[
 (\omega^s)^{-1}-\widetilde\omega(s,0)
 =\frac{s^2}{G_{12}(s,0)G_{22}(s,0)},
\]
whose absolute value is at most $Cs^2/r_s^4\le Cs/r_s^2$.
The analogous segment and twist estimates therefore give
\begin{equation}\label{notes-twist-match-second}
 |\widetilde\omega-(\omega^s)^{-1}|\le Cs/r_s^2
 \quad\text{on }\mathcal B_s^+.
\end{equation}

Now we can construct the adapted Killing fields. For a constant vector $\mathbf{v}=(-c_0,1)^T$, set $K=-c_0X_1+X_2 = (X_1, X_2) \mathbf{v}$. Since 
$$d \theta ^{1}(K) = -c_0, d \theta ^2 (K) = 1, (d \theta ^1 + \omega d \theta ^2)(K) = \omega - c_0,$$ 
by \eqref{notes-axis-gram},
 $$|K|_g^2=k \rho ^2(\omega -c_0)^2+ k ^{-1} .$$
Also,
$$K ^{\flat_g} = (G \mathbf{v})_i d \theta ^{i} = k\rho^2 (\omega - c_0) (d \theta ^{1} + \omega d \theta ^2) + k ^{-1} d \theta ^2 .$$
So
\begin{align*}
dK^{\flat_g}
&=d\left(k\rho^2(\omega-c_0)\right)
       \wedge(d\theta^1+\omega d\theta^2)
       +k\rho^2(\omega-c_0)d\omega\wedge d\theta^2
-k^{-2}dk\wedge d\theta^2\\
   &=\left( \rho^2(\omega-c_0)dk +2k\rho(\omega-c_0)d\rho +k\rho^2 d\omega \right) 
   \wedge(d\theta^1+\omega d\theta^2)\\
&\ \ +\left( k\rho^2(\omega-c_0)d\omega-k^{-2}dk\right) 
\wedge d\theta^2.
\end{align*} 
Since $K$ is Killing, 
\begin{align*}
	2|\nabla ^{g}K|^2_{g} &=  |d K ^{\flat_g}|^2_g \\
			  &= \frac{1}{k \rho ^2}\left| \left( \rho^2(\omega-c_0)dk
+2k\rho(\omega-c_0)d\rho
+k\rho^2 d\omega \right)\right|^2_h \\
			  &\ \ + k \left| \left( 
k\rho^2(\omega-c_0)d\omega-k^{-2}dk\right)\right|^2_h.
\end{align*} 
Consequently,
\begin{align*}
	|\nabla ^{g} K|^2_g &\leq C e ^{- 2 \nu } \left( k\rho ^2(\omega -c_0)^2 |d \log k|^2_{\delta _2} + k (\omega -c_0)^2 |d \rho |^2_{\delta _2} + k \rho ^2 |d \omega |^2_{\delta _2} \right) \\
			    &\ \ + C e ^{- 2 \nu } \left( k^3 \rho ^{4} (\omega -c_0)^2|d \omega |^2_{\delta _2} + k ^{-1} |d \log k|^2_{\delta _2} \right) \\
			    &\leq C e ^{-2 \nu }\left(  |K|^2_g |d \log k|^2_{\delta _2}+ k ^2 \rho ^2 |K|^2_g |d \omega |^2_{\delta _2} + k(\omega -c_0)^2 \right) .
\end{align*} 
Thus
\begin{align}\label{notes-K-ratio}
	\begin{split}
		\frac{|\nabla ^{g}K|_g^2}{|K|_g^2}
		& \leq C e ^{-2 \nu } \left( |d \log k|^2_{\delta _2} + k ^2 \rho ^2 |d \omega |^2_{\delta _2} + k^2(\omega -c_0)^2 \right) .
	\end{split}
\end{align}

On $\mathcal B_s^-$ take $c_0=\omega^s$.
Since $k\le Cr_s^2/s^2$, \eqref{notes-k-control},
\eqref{notes-twist-gradient}, and \eqref{notes-twist-match} bound each
term in parentheses by $C/s^2$. Thus \eqref{notes-annular-conformal}
gives $|\nabla^gK_s|_g/|K_s|_g\le C/r_s$.
On $\mathcal B_s^+$ write
$K_s=-\omega^s(X_1-(\omega^s)^{-1}X_2)$ and use
\eqref{notes-twist-match-second} in the same way.
These two bands cover the part of $s\le\xi\le2s$ and $\rho<s/2$. 


It remains to control the part $s\le\xi\le2s$ and $\rho\ge s/2$. Note that 
$$dK^{\flat_g}
=\sum_{\alpha=\rho,\tau}\sum_i(\partial _\alpha G\mathbf v)_i\,dx^\alpha\wedge d\theta^i,$$
so
\begin{align*}
	|\nabla ^{g}K|^2_g &= \frac12e^{-2\nu}
\sum_{\alpha=\rho,\tau}
(\partial _\alpha G\mathbf v)^TG^{-1}(\partial _\alpha G\mathbf v) \\
			   & \leq \frac{1}{2} e ^{-2 \nu } \sum_{\alpha }^{} \mathrm{tr}(G ^{-1}\partial _\alpha G G ^{-1} \partial _\alpha G) \mathbf{v}^{T}G \mathbf{v}\\
			   &= \frac{1}{2} e ^{-2 \nu } \sum_{\alpha }^{} \left( 2\rho ^{-2} (\partial _\alpha \rho )^2 + 2|H_{\alpha }|^2_{\mathbb{H}^2} \right)|K|^2_g\\
			   &\leq e ^{-2 \nu }(\rho ^{-2} + |d H|^2_{\delta _3, \mathbb{H}^2})|K|^2_g.
\end{align*} 
The Weyl equation gives
$|dH|_{\delta_3,\mathbb H^2}\le C/\rho $ when $s \leq \xi \leq 2s$. If $\rho \geq s/ 2$, we also have $|\nabla ^{g}K_s|_g /|K_s|_g \leq C e ^{-\nu } /s \leq C / r_s$.
This proves the stated estimate on all of $\{s\le\tilde\xi\le2s\}$.
\end{proof}

\subsection{Curvature energy and the classification}
We apply the local Euler and Pontryagin transgression formulas for a
Killing field \cite[Sections~2.1--2.2]{notes-Weber}. The Euler identity first
establishes finite curvature energy; the Pontryagin identity then forces
one Weyl half to vanish.

Let $\mathcal A_s=\{s\le\tilde\xi\le2s\}\subset M$, and recall $r_s=se^{\nu(s,0)}$.
So
\begin{align}\label{notes-annular-collapse}
	\begin{split}
		\Vol_g(\mathcal A_s) &= 4 \pi ^2 \int_{\substack{s\le\xi\le2s\\\rho>0}} e^{2 \nu } \rho d \rho d \tau \leq C e ^{2\nu (s,0)} s ^3 \le C F(s, 0) s^2,\\
		v_s&:=r_s^{-4}\Vol_g(\mathcal A_s)
       \le\frac{C}{F(s,0)}\to 0.
	\end{split} 
\end{align}

For a Ricci-flat four-metric the Euler form is
\[
 e_g=\frac1{8\pi^2}
       (|W^+|_{g,\mathrm{op}}^2+|W^-|_{g,\mathrm{op}}^2)\,dV_g\ge0.
\]
Here the norms are Hilbert--Schmidt norms of the curvature operators on
$\Lambda_g^\pm$; the four-index tensor norm satisfies
$|\Rm_g|_g^2=4(|W^+|_{g,\mathrm{op}}^2+|W^-|_{g,\mathrm{op}}^2)$.

\begin{prop}\label{notes-Euler-energy}
Assume $g$ is not flat. Then
\[
 \int_M e_g=1,\qquad \int_M|\Rm_g|_g^2\,dV_g = 32 \pi ^2<\infty.
\]
\end{prop}
\begin{proof}
	Let $K_s= -\omega^s X_1+ X_2$ be the Killing field in previous lemma. 
Both coefficients of $K_s$ are nonzero, so its only zero is $p$.
Its linearization there is a pair of nonzero planar rotations and has
positive determinant, hence index one. More precisely, using the torus action splitting of $T_pM$, $(\nabla ^{g}_{\cdot } K_s)_p: T_pM \to T_p M$ has the form
\begin{align*}
	\begin{pmatrix} -\omega^s J & 0 \\ 0 & J \end{pmatrix} , \ \ J = \begin{pmatrix} 0 & -1 \\ 1 & 0 \end{pmatrix} ,
\end{align*}
so $\mathrm{det}(\nabla ^{g}_{\cdot } K_s)_p = (\omega^s)^2 >0	.$
We use Chern's classical Euler transgression
\cite[Sections~2--3]{notes-Chern1944},
\cite[Section~1, (9)--(11), and Section~2]{notes-Chern1945}.
Let $\varpi:S(TM)\to M$ be the unit tangent sphere bundle. Chern's angular
three-form $\Psi$ satisfies
\[
 d\Psi=-\varpi^*e_g,\qquad \int_{S(T_xM)}\Psi=1,
\]
where each fiber has its boundary orientation as the unit sphere in $T_xM$.
The unit vector field $u_s=K_s/|K_s|_g$ is a section of this bundle on
$M\setminus\{p\}$. Therefore the three-form
\[
 T_s^e:=-u_s^*\Psi
 \quad\hbox{satisfies}\quad dT_s^e=e_g
 \quad\hbox{on }M\setminus\{p\}.
\]
In dimension four, Chern's formula contains terms cubic in $\nabla^g u_s$
and terms linear in both $\nabla^g u_s$ and curvature. Since
$|\nabla^g u_s|_g\le |\nabla^gK_s|_g/|K_s|_g$, it gives
\[
 |T_s^e|_g\le C(\ell_s|\Rm_g|_g+\ell_s^3),\qquad
 \ell_s:=|\nabla^gK_s|_g/|K_s|_g.
\]
Nondegeneracy of the zero gives $\ell_s=O(r^{-1})$ for
$r=d_g(p,\cdot)$, so $T_s^e$ is locally integrable at $p$.
The fiber normalization and the local index give
\[
 \lim_{\epsilon\downarrow0}\int_{\partial B_g(p,\epsilon)}T_s^e
 =-\operatorname{ind}_p(K_s)=-1,
\]
with the sphere oriented outward from $p$. Stokes's theorem on the
complement of this ball consequently yields
\[
 dT_s^e=e_g-\delta_p
 \quad\hbox{distributionally on }M,
\]
where $\delta_p$ is the Dirac four-current, $\langle\delta_p,\varphi\rangle=\varphi(p)$.

Choose $\eta_s=\eta(\xi/s)$, with $\eta=1$ on $[0,1]$, $\eta=0$ on
$[2,\infty)$ and $0\le\eta\le1$. On $\mathcal A_s$, 
$$|d \eta _s| = s ^{-1}|\eta '| \cdot |d \xi |_g \leq C s ^{-1} e ^{- \nu } \leq C s ^{-1} e ^{- \nu (s,0)} = C r_s ^{-1},$$
so $|d\eta_s|_g+\ell_s\le Cr_s ^{-1}$.
Set $I_s=\int_M\eta_s^2e_g$. Since $\eta _s(p) =1$,
\begin{align*}
	I_s - 1 &= - \int_{M} d \eta _s ^2 \wedge T ^{e}_{s} = -2 \int_{\mathcal{A}_s} \eta _s d \eta _s \wedge T^e_{s}.
\end{align*} 
Cauchy--Schwarz and \eqref{notes-annular-collapse} give
\begin{align*}
 |I_s-1|
 &\le Cr_s^{-2}\int_{\mathcal A_s}\eta_s|\Rm_g|_g\,dV_g+Cv_s\\
 &\le C\sqrt{v_s I_s}+Cv_s.
\end{align*}
As $s\to \infty$,
Young's inequality first bounds $I_s$, and the same estimate then gives
$I_s\to1$. Properness of $\xi$ gives $\eta_s\to1$ pointwise. Fatou's lemma
establishes finite total Euler energy, after which dominated convergence
shows that its value is one.
\end{proof}

Define the two nonnegative chiral energies
\[
 \mathcal E_\pm=\frac1{8\pi^2}\int_M|W^\pm|_{g,\mathrm{op}}^2dV_g,
 \qquad \mathcal E_++\mathcal E_-=1.
\]
For the curvature two-form $\Omega^g$ of $\nabla^g$, our first Pontryagin form is
\[
 p_1(\nabla^g)=-\frac1{8\pi^2}\tr(\Omega^g\wedge\Omega^g),\qquad
 P:=\int_Mp_1=2(\mathcal E_+-\mathcal E_-).
\]
In particular, $|P|\le2$, with equality exactly when one Weyl half vanishes.

\begin{proof}[Proof of Theorem~\ref{notes-R4-classification}]
It remains to consider the nonflat case. As in the Euler argument, we
transgress the curvature form away from the zero and compute its local
residue. For the Pontryagin form we use the Killing-field construction in
\cite[Section~2.1, (3), p.~3119]{notes-Weber}.
Let $K=a_1X_1+a_2X_2$, where $a_1a_2\ne0$, and on $M\setminus\{p\}$ set
\[
 \alpha=\frac{K^{\flat_g}}{|K|_g^2},\qquad B(V)=\nabla_V^gK.
\]
Here $B$ is a skew-adjoint endomorphism field. The transgression three-form is
\begin{equation}\label{notes-P-form}
 T_K=-\frac1{4\pi^2}\alpha\wedge\tr(B\Omega^g)
       +\frac1{8\pi^2}\tr(B^2)\alpha\wedge d\alpha,
\end{equation}
where the trace is over tangent-space endomorphisms.
To see why this form exists, compare $\nabla^g$ with the metric connection
\[
 \nabla^0_VY=\nabla^g_VY-\alpha(V)B(Y).
\]
Choose locally an orthonormal frame invariant under the flow of $K$.
If $\vartheta$ is the connection matrix of $\nabla^g$ in this frame, then
$\vartheta(K)=B$, because $[K,e_i]=0$ implies
$\nabla_K^g e_i=\nabla_{e_i}^gK$. Hence the connection matrix
$\vartheta^0=\vartheta-\alpha B$ is invariant and annihilates $K$.
Cartan's formula gives $\iota_K\Omega^0=0$ for its curvature $\Omega^0$.
Consequently $\tr(\Omega^0\wedge\Omega^0)=0$ in dimension four.
The standard transgression between these two connections gives
\eqref{notes-P-form} and $dT_K=p_1(\nabla^g)$ on $M\setminus\{p\}$.
Moreover,
\[
 |\alpha|_g=|K|_g^{-1},\qquad
 |d\alpha|_g\le C\frac{|\nabla^gK|_g}{|K|_g^2},
\]
so \eqref{notes-P-form} gives
\[
 |T_K|_g\le C(\ell_K|\Rm_g|_g+\ell_K^3),\qquad
 \ell_K=\frac{|\nabla^gK|_g}{|K|_g}.
\]

Use the orientation of the ordered primitive normal-coordinate planes at
$p$, with signed rotation speeds $a_1,a_2$. In their tangent model
$(\mathbb C^2,\delta_4)$, where $\delta_4$ is Euclidean, write
$K_0(z_1,z_2)=(ia_1z_1,ia_2z_2)$ and
$\alpha_0=K_0^{\flat_{\delta_4}}/|K_0|_{\delta_4}^2$. The linearization
$B_0=\operatorname{diag}(a_1J,a_2J)$ satisfies
$\tr(B_0^2)=-2(a_1^2+a_2^2)$. On the unit sphere, putting
$u=|z_1|^2$ gives
\[
 \int_{S^3}\alpha_0\wedge d\alpha_0
 =4\pi^2a_1a_2\int_0^1
       \frac{du}{[a_1^2u+a_2^2(1-u)]^2}
 =\frac{4\pi^2}{a_1a_2}.
\]
The spheres are oriented outward. In equivariant normal coordinates,
$g=\delta_4+O(r^2)$ and $B=B_0+O(r^2)$, where $r=d_g(p,\cdot)$.
The curvature term in \eqref{notes-P-form} has integral $O(r^2)$ on the
sphere of radius $r$, and the remaining term converges to its Euclidean
model. Thus
\[
 \lim_{\epsilon\downarrow0}
 \int_{\partial B_g(p,\epsilon)}T_K
 =\frac{\tr(B_0^2)}{8\pi^2}\int_{S^3}\alpha_0\wedge d\alpha_0
 =-\left(\frac{a_1}{a_2}+\frac{a_2}{a_1}\right).
\]
Since $\ell_K=O(r^{-1})$, the form $T_K$ is locally integrable at $p$.
Stokes's theorem, as in the Euler case, gives
\begin{equation}\label{notes-P-transgression}
 dT_K=p_1(\nabla^g)
       -\left(\frac{a_1}{a_2}+\frac{a_2}{a_1}\right)\delta_p
 \quad\hbox{distributionally on }M.
\end{equation}

Apply \eqref{notes-P-transgression} to $K_s=-\omega^sX_1+X_2$ and a cutoff
$\chi_s$ equal to one on $\xi\le s$, zero on $\xi\ge2s$, with
$|d\chi_s|_g\le C/r_s$. The annular Killing estimate gives
\begin{align*}
 \left|\int_M\chi_s p_1+\omega^s+(\omega^s)^{-1}\right|
 &\le C\left[
 r_s^{-2}\Vol_g(\mathcal A_s)^{1/2}
       \left(\int_{\mathcal A_s}|\Rm_g|_g^2dV_g\right)^{1/2}
       +v_s\right]\\
 &\to 0.
\end{align*}
Here \eqref{notes-annular-collapse} and Proposition~\ref{notes-Euler-energy}
justify the limit. Absolute integrability of $p_1$ now yields
\[
 -\omega^s-(\omega^s)^{-1}\to  P.
\]
For every nonzero real $t$, $|t+t^{-1}|\ge2$. Thus $|P|\ge2$, while
$\mathcal E_++\mathcal E_-=1$ gave $|P|\le2$. Equality forces
$W^+=0$ or $W^-=0$.

After choosing orientation, the induced connection on $\Lambda_g^+$ is
flat. Since $M$ is simply connected, it has a global parallel orthonormal
frame of self-dual forms. Scaling these forms to norm $\sqrt2$ gives a
hyperk\"ahler triple. The remaining classification is a consequence of the
known complete hyperk\"ahler circle classification: the isometric torus
supplies the required tri-Hamiltonian circle as follows. The torus acts on
this three-dimensional space through $SO(3)$. An abelian connected subgroup
of $SO(3)$ has dimension at most one, so the kernel contains a closed circle
preserving the triple. For its generator $K$, the one-forms
$\iota_K\omega_i$ are closed by Cartan's formula and exact because
$H^1(M;\mathbb R)=0$. Thus this is a tri-Hamiltonian circle action.

The complete hyperk\"ahler circle classification of Swann
\cite[Theorem~3.1]{notes-Swann}, with the center description in
Dancer--Swann \cite[Section~2.2]{notes-DancerSwann}, applies to this complete tri-Hamiltonian circle action. In this description the empty
center case has topology $S^1\times\mathbb R^3$, and a nonempty center set
$Z$ has $b_2=|Z|-1$ (infinitely many centers give infinite second Betti
number). Since $M$ is homeomorphic to $\mathbb R^4$, there is exactly one
center. With the circle generator normalized to period $2\pi$, the
one-center Gibbons--Hawking potential is
\[
 V(x)=c_0+\frac1{2|x-x_0|},\qquad c_0\ge0,
\]
where $x$ is a moment-map coordinate on $\mathbb R^3$. The case $c_0=0$
is Euclidean and $c_0>0$ is Taub--NUT; all positive one-center parameters
differ by constant scaling.
The nonflat branch is therefore Taub--NUT.
\end{proof}

\begin{thebibliography}{99}

\bibitem{notes-AKP}
S. Alexander, V. Kapovitch and A. Petrunin,
\emph{Invitation to Alexandrov geometry: CAT[0] spaces}.

\bibitem{notes-Besse}
A. L. Besse, \emph{Einstein Manifolds},
Ergebnisse der Mathematik und ihrer Grenzgebiete (3), vol.~10,
Springer-Verlag, Berlin, 1987.

\bibitem{notes-BH}
M. R. Bridson and A. Haefliger,
\emph{Metric Spaces of Non-Positive Curvature}, Springer, 1999.

\bibitem{notes-CheegerTian}
J. Cheeger and G. Tian,
\emph{Curvature and injectivity radius estimates for Einstein 4-manifolds},
J. Amer. Math. Soc. \textbf{19} (2006), no.~2, 487--525.
%\href{https://doi.org/10.1090/S0894-0347-05-00511-4}{doi:10.1090/S0894-0347-05-00511-4}.

\bibitem{notes-ChenLi}
X. Chen and Y. Li, \emph{On the geometry of asymptotically flat manifolds},
Geom. Topol. \textbf{25} (2021), no.~5, 2469--2572.
%\href{https://doi.org/10.2140/gt.2021.25.2469}{doi:10.2140/gt.2021.25.2469}.
%\href{https://arxiv.org/abs/1908.07248v3}{arXiv:1908.07248v3}.

\bibitem{notes-Cheng}
S.-Y. Cheng, \emph{Eigenvalue comparison theorems and its geometric applications},
Math. Z. \textbf{143} (1975), 289--297.
%\href{https://doi.org/10.1007/BF01214381}{doi:10.1007/BF01214381}.

\bibitem{notes-Chern1944}
S.-S. Chern, \emph{A simple intrinsic proof of the Gauss--Bonnet formula
for closed Riemannian manifolds},
Ann. of Math. (2) \textbf{45} (1944), no.~4, 747--752.
%\href{https://doi.org/10.2307/1969302}{doi:10.2307/1969302}.

\bibitem{notes-Chern1945}
S.-S. Chern, \emph{On the curvatura integra in a Riemannian manifold},
Ann. of Math. (2) \textbf{46} (1945), no.~4, 674--684.
%\href{https://doi.org/10.2307/1969203}{doi:10.2307/1969203}.

\bibitem{notes-Colding}
T. H. Colding, \emph{New monotonicity formulas for Ricci curvature and applications. I},
Acta Math. \textbf{209} (2012), 229--263.
%\href{https://doi.org/10.1007/s11511-012-0086-2}{doi:10.1007/s11511-012-0086-2}.

\bibitem{notes-DancerSwann}
A. Dancer and A. Swann, \emph{Hypertoric manifolds and hyperK\"ahler moment maps},
in \emph{Special Metrics and Group Actions in Geometry},
Springer INdAM Series, vol.~23, Springer, Cham, 2017, 107--127.
%\href{https://doi.org/10.1007/978-3-319-67519-0_5}{doi:10.1007/978-3-319-67519-0\_5}.
%\href{https://arxiv.org/abs/1607.04078v1}{arXiv:1607.04078v1};

\bibitem{notes-GibbonsHawking}
G. W. Gibbons and S. W. Hawking, \emph{Gravitational multi-instantons},
Phys. Lett. B \textbf{78} (1978), no.~4, 430--432.
%\href{https://doi.org/10.1016/0370-2693(78)90478-1}{doi:10.1016/0370-2693(78)90478-1}.

\bibitem{notes-Hatcher}
A. Hatcher, \emph{The Kirby torus trick for surfaces},
\href{https://arxiv.org/abs/1312.3518v3}{arXiv:1312.3518v3}.

\bibitem{notes-Hawking}
S. W. Hawking, \emph{Gravitational instantons},
Phys. Lett. A \textbf{60} (1977), no.~2, 81--83.
%\href{https://doi.org/10.1016/0375-9601(77)90386-3}{doi:10.1016/0375-9601(77)90386-3}.

\bibitem{notes-LiSunCompact}
M. Li and S. Sun, \emph{Compact Toric Einstein Four-Manifolds},
manuscript.

\bibitem{notes-LiSunHarmonic}
M. Li and S. Sun, \emph{Gravitational instantons and harmonic maps},
\href{https://arxiv.org/abs/2507.15284v1}{arXiv:2507.15284v1}, 2025.

\bibitem{notes-LiYau}
P. Li and S.-T. Yau, \emph{On the parabolic kernel of the Schr\"odinger operator},
Acta Math. \textbf{156} (1986), 153--201.
%\href{https://doi.org/10.1007/BF02399203}{doi:10.1007/BF02399203}.

\bibitem{notes-Liu}
T. Liu, \emph{Toric Einstein 4-manifolds with non-negative sectional curvature},
Ann. Global Anal. Geom. \textbf{67} (2025), article~14.
%\href{https://doi.org/10.1007/s10455-025-09990-3}{doi:10.1007/s10455-025-09990-3}.

\bibitem{notes-OrlikRaymond}
P. Orlik and F. Raymond, \emph{Actions of the torus on 4-manifolds. I},
Trans. Amer. Math. Soc. \textbf{152} (1970), 531--559.

\bibitem{notes-Swann}
A. Swann, \emph{Twists versus modifications},
Adv. Math. \textbf{303} (2016), 611--637.
%\href{https://doi.org/10.1016/j.aim.2016.08.028}{doi:10.1016/j.aim.2016.08.028}.
%\href{https://arxiv.org/abs/1512.06343v1}{arXiv:1512.06343v1};

\bibitem{notes-Weber}
B. Weber, \emph{Energy and asymptotics of Ricci-flat 4-manifolds with a Killing field},
Proc. Amer. Math. Soc. \textbf{147} (2019), 3117--3130.
%\href{https://doi.org/10.1090/proc/14014}{doi:10.1090/proc/14014}.
%\href{https://arxiv.org/abs/1308.4997v2}{arXiv:1308.4997v2}.

\bibitem{notes-Yau}
S.-T. Yau, \emph{Harmonic functions on complete Riemannian manifolds},
Comm. Pure Appl. Math. \textbf{28} (1975), 201--228.
%\href{https://doi.org/10.1002/cpa.3160280203}{doi:10.1002/cpa.3160280203}.

\end{thebibliography}
\end{document}
